\documentclass[11pt]{article}

\usepackage[T1]{fontenc}
\usepackage{lmodern}
\usepackage{microtype}
\usepackage[a4paper,margin=27mm]{geometry}
\emergencystretch=3em
\usepackage{amsmath,amssymb,amsthm,mathtools,mathrsfs}
\usepackage{enumitem}
\usepackage{booktabs,longtable,array}
\usepackage{xcolor,xurl}
\usepackage{aliascnt}
\usepackage{hyperref}
\usepackage[nameinlink,capitalise,noabbrev]{cleveref}

\hypersetup{
 colorlinks=true,
 linkcolor=blue!55!black,
 citecolor=blue!55!black,
 urlcolor=blue!55!black,
 pdftitle={Multiweight arithmetic holonomy and density-one irrationality of p-adic zeta values},
 pdfauthor={Christopher D. Long}
}

\newtheorem{theorem}{Theorem}[section]
\newaliascnt{proposition}{theorem}
\newtheorem{proposition}[proposition]{Proposition}
\aliascntresetthe{proposition}
\newaliascnt{lemma}{theorem}
\newtheorem{lemma}[lemma]{Lemma}
\aliascntresetthe{lemma}
\newaliascnt{corollary}{theorem}
\newtheorem{corollary}[corollary]{Corollary}
\aliascntresetthe{corollary}
\newaliascnt{claim}{theorem}
\newtheorem{claim}[claim]{Claim}
\aliascntresetthe{claim}
\theoremstyle{definition}
\newaliascnt{definition}{theorem}
\newtheorem{definition}[definition]{Definition}
\aliascntresetthe{definition}
\theoremstyle{remark}
\newaliascnt{remark}{theorem}
\newtheorem{remark}[remark]{Remark}
\aliascntresetthe{remark}
\newaliascnt{warning}{theorem}
\newtheorem{warning}[warning]{Warning}
\aliascntresetthe{warning}

\crefname{theorem}{Theorem}{Theorems}
\crefname{proposition}{Proposition}{Propositions}
\crefname{lemma}{Lemma}{Lemmas}
\crefname{corollary}{Corollary}{Corollaries}
\crefname{claim}{Claim}{Claims}
\crefname{definition}{Definition}{Definitions}
\crefname{remark}{Remark}{Remarks}
\crefname{warning}{Warning}{Warnings}

\newcommand{\Q}{\mathbf Q}
\newcommand{\Z}{\mathbf Z}
\newcommand{\C}{\mathbf C}
\newcommand{\R}{\mathbf R}
\newcommand{\F}{\mathbf F}
\newcommand{\Qp}{\mathbf Q_p}
\newcommand{\dd}{\,\mathrm d}
\newcommand{\ord}{\operatorname{ord}}
\newcommand{\rank}{\operatorname{rank}}
\newcommand{\lcmop}{\operatorname{lcm}}
\newcommand{\Ann}{\operatorname{Ann}}
\newcommand{\Sym}{\operatorname{Sym}}
\newcommand{\Span}{\operatorname{Span}}
\newcommand{\calE}{\mathcal E}
\newcommand{\calF}{\mathcal F}
\newcommand{\calG}{\mathcal G}
\newcommand{\calW}{\mathcal W}
\newcommand{\calH}{\mathcal H}
\newcommand{\calA}{\mathcal A}
\newcommand{\calM}{\mathcal M}
\newcommand{\calO}{\mathcal O}
\newcommand{\calV}{\mathcal V}
\newcommand{\widehatdeg}{\widehat{\deg}}
\newcommand{\Energy}{\mathscr H}
\newcommand{\LambdaAH}{\Lambda}
\newcommand{\floor}[1]{\left\lfloor #1\right\rfloor}
\newcommand{\ceil}[1]{\left\lceil #1\right\rceil}
\newcommand{\abs}[1]{\left|#1\right|}
\newcommand{\positivepart}[1]{\left(#1\right)_+}
\newcommand{\T}{\mathbb T}

\title{Multiweight arithmetic holonomy and\\
density-one irrationality of $p$-adic zeta values}
\author{Christopher D. Long}
\date{All-prime dependency-audited exact-capacity revision --- 26 August 2026}

\begin{document}
\maketitle

\begin{abstract}
For each genus-zero prime
\[
 \mathcal P_0=\{2,3,5,7,13\}
\]
and each finite set $\mathcal S$ of distinct odd integers at least $3$, we
combine the Eisenstein--Eichler extensions attached to the values
$\zeta_p(s)$, $s\in\mathcal S$, using one common homogenizing coordinate.
Under the joint rationality hypothesis the resulting rank-$m$ horizontal
leaf, where $m=1+\sum_{s\in\mathcal S}s$, has full functional rank.  The
blockwise genuine-source Hasse kernels and a common torsor-valued Cartier
argument give the Hasse-improved multiweight type
\[
 \tau^{\mathrm H}_{p,\mathcal S}(\xi)
 =\frac{2}{m^2}\left[
 \xi\sum_{s\in\mathcal S}s^2
 -\left(\sum_{s\in\mathcal S}(s-1)\right)J_p(\xi)
 +(\max\mathcal S)I^m_\xi(\xi)
 \right].
\]
A thickness-sensitive Fitting refinement replaces the rectangular pole bound
by a denominator polygon.  The certified finite conclusions are unchanged:
among the first $1200$ positive odd arguments, at least
\[
1194,\quad1191,\quad1185,\quad1180,\quad1061
\]
of the values $\zeta_p(s)$ are irrational for
$p=2,3,5,7,13$, respectively.

We then develop an intrinsic-curve version that no longer assumes
$X_0(p)\simeq\mathbf P^1$.  On the fixed coarse modular curve, after
removing the finite elliptic locus, a finite-map restriction-of-scalars
argument gives a Shidlovsky zero estimate and exact
$q\mapsto q^\ell$ divided Frobenius gives a general-curve $q$-jet Cartier
window.  A normalization-robust slopes argument applies finite-order
capacitary estimates only after contraction to a scalar section.  It yields
\[
 \tau^{\mathrm{curve}}_{p,\mathcal S}(\xi)
 \ge \Lambda_p^\circ-\frac{C_p}{m},
\]
where $\Lambda_p^\circ>0$ comes from the strict capacitary gain of a
$p$-adic annular collar and $C_p$ is independent of the packet.

We also construct the canonical exact metric.  Its finite components are the
Bost--Chambert-Loir capacitary $\Q$-model metrics, while its Archimedean
components are the Bost--Charles direct-image Green functions of the chosen
continuation maps.  After clearing the finitely many rational exponents, the
capacitary tangent line is the normal line of a pseudoconcave formal-analytic
surface.  The Chen--Moriwaki determinant Hilbert--Samuel theorem and a
uniform finite-order Green--Jensen estimate then give the exact source formula
and the target-normal cancellation
\[
 \widehatdeg M_D=\frac m2\Energy_\alpha D^2+o(D^2),
 \qquad
 \tau^{\mathrm{curve}}_{p,\mathcal S}(\xi)
 \ge \LambdaAH_\alpha-\frac{\Energy_\alpha}{m}.
\]
The stronger exact theorem is not needed for the normalization-robust
qualitative proof.  We then compute the maximal modular capacity pair exactly
at every prime.  For $0<Y<1/p$,
\[
 (\Lambda_p(Y),\Energy_p(Y))=
 \left(\frac{12\log p}{p-1}-2\pi Y,
       \frac{12\log p}{p-1}-2\pi Y+C_p(Y)\right).
\]
For $p\ge5$ the finite contribution is the effective resistance of the
supersingular reduction graph, evaluated by the Eichler--Deuring mass formula;
$p=2,3$ are treated directly with their integral Hauptmodul lemniscates.  If
$d_p$ is the degree of the finite map used in the general-curve zero estimate,
then
\[
 R_p(X):=\#\{3\le s\le X:\ s\text{ odd and }\zeta_p(s)\in\mathbf Q\}
 \le\left(\frac{e\,d_p(p-1)}{12\log p}+o_p(1)\right)(\log X)^2.
\]
Thus the irrational values have natural density one with an explicit leading
constant at every fixed prime.  A fixed-point finite-packet certificate gives,
among the first $1200$ odd arguments, at least
\[
 1010,974,963,823,786,775,743,605,594
\]
irrational values for $p=11,17,19,23,29,31,37,41,43$, respectively.
\end{abstract}

\medskip
\noindent\textbf{Audit status.}
The genus-zero part isolates the same published interfaces as the underlying
single-weight arithmetic-holonomy construction: Calegari's Eisenstein family,
Buzzard continuation, Katz's Hasse invariant, the de Rham frame torsor of
Graham--Pilloni--Rodrigues Jacinto, and the Bost--Charles/CDT slopes and
prime-window estimates.  The all-prime extension uses Bost--Chambert-Loir
capacitary metrics and four internal bridges: coarse-moduli descent and
finite-map Shidlovsky restriction of scalars, intrinsic $q$-jet Cartier
strictness, scalar finite-order capacitary jet bounds with a fixed-twist
source comparison, and strict positivity of the distinguished $p$-adic
capacitary gain.  The compatibility of these bounds with the hybrid
finite-prime lattices is proved explicitly.

The exact normalization is proved separately in
\cref{thm:curve-Bost-exact,app:exact-normalization}.  The audit distinguishes
two assertions that had previously been conflated.  Bost--Chambert-Loir
provides one adelic Arakelov $\Q$-divisor inducing all finite capacitary
metrics; after denominator clearing, its metrized tangent line is realizable
as the normal line of a pseudoconcave Bost--Charles formal-analytic surface.
We do not assert that the entire finite vertical divisor is literally a
Bost--Charles direct image, nor that matching the normal line automatically
produces a Bost--Charles meromorphic map to the modular curve.  Instead, the
exact source and jet formulas are proved directly on the modular curve for the
canonical mixed capacitary/direct-image divisor, using the Chen--Moriwaki
determinant Hilbert--Samuel theorem and a place-by-place finite-order
Green--Jensen estimate.  The exact modular capacity pair is computed for every prime in
\cref{thm:allprime-capacity-pair}.  Its maximal-domain passage uses the
rational skeleton exhaustion of \cref{lem:maximal-capacity-exhaustion};
Bost--Chambert-Loir Proposition~5.6 is used only to realize each fixed inner
domain as a rational lemniscate, not as a continuity theorem.  Supersingular
residue fields are split by normalized unramified base change.  Consequently
the all-prime density coefficient is explicit in terms of the finite-map
degree $d_p$.  The positive-genus finite-packet claims use the certified upper
bound $2g_p\ge d_p$ and are independently fixed-point certified.  Independent
specialist review remains appropriate before publication.
\tableofcontents

\section{Introduction and main results}\label{sec:intro}

For an odd integer $s\ge3$, let $\zeta_p(s)$ denote the Kubota--Leopoldt
$p$-adic zeta value in Calegari's normalization.  The single-weight
Eisenstein--Eichler method starts from
\[
 K_{p,s}(q)=\frac{\zeta_p(s)}2+
 \sum_{n\ge1}\sigma_{-s}^{(p)}(n)q^n,
 \qquad
 \sigma_{-s}^{(p)}(n)=\sum_{\substack{d\mid n\\p\nmid d}}d^{-s}.
\]
Under the contrary hypothesis $\zeta_p(s)\in\Q$, this is an arithmetic
horizontal leaf of a rank-$(s+1)$ algebraic connection.  The present paper
uses several such blocks simultaneously but shares their scalar coordinate.
The resulting rank is essentially the sum of the weights, while the largest
large-prime pole grade is only their maximum.

Let $\mathcal S$ be a finite nonempty set of distinct odd integers $s\ge3$.
Put
\begin{equation}\label{eq:packet-data}
 m=1+\sum_{s\in\mathcal S}s,
 \qquad Q_{\mathcal S}=\sum_{s\in\mathcal S}s^2,
 \qquad B_{\mathcal S}=\sum_{s\in\mathcal S}(s-1),
 \qquad s_*=\max\mathcal S.
\end{equation}
For $p\in\{2,3,5,7,13\}$, define
\begin{equation}\label{eq:Jp}
 c_p=\frac{p+1}{12},
 \qquad
 J_p(\xi)=\int_0^\xi(1-c_px)_+\,\dd x.
\end{equation}
Thus
\[
 J_p(\xi)=
 \begin{cases}
 \xi-\dfrac{p+1}{24}\xi^2,&c_p\xi\le1,\\[2mm]
 \dfrac{6}{p+1},&c_p\xi\ge1.
 \end{cases}
\]
In particular, because $c_{13}=7/6$, every admissible cutoff $\xi>1$ satisfies
\begin{equation}\label{eq:J13}
 J_{13}(\xi)=\frac37.
\end{equation}

For $m\ge2$ and $1<\xi<m$, put
\begin{equation}\label{eq:I-def}
 K_m(\xi)=\floor{\frac{m-1}{\xi}},
\end{equation}
and
\begin{equation}\label{eq:I}
 I^m_\xi(\xi)=
 \left(m-\frac12\right)H_{K_m(\xi)}-K_m(\xi)\xi
 +\frac{\positivepart{m-(K_m(\xi)+1)\xi}^2}{2(K_m(\xi)+1)}.
\end{equation}

\begin{theorem}[Hasse-improved multiweight arithmetic holonomy]
\label{thm:multiweight}
Fix $p\in\{2,3,5,7,13\}$ and a finite set $\mathcal S$ as above.  Assume
\[
 \zeta_p(s)\in\Q\qquad(s\in\mathcal S).
\]
Let $\boldsymbol\phi$ be admissible adelic continuation templates for the
simultaneous Eisenstein--Eichler leaf, with derivative width
$\LambdaAH(\boldsymbol\phi)$ and Bost--Charles energy
$\Energy(\boldsymbol\phi)$.  For every $1<\xi<m$,
\begin{equation}\label{eq:multiweight-holonomy}
 m\left(\LambdaAH(\boldsymbol\phi)
 -\tau^{\mathrm H}_{p,\mathcal S}(\xi)\right)
 \le\Energy(\boldsymbol\phi),
\end{equation}
where
\begin{equation}\label{eq:tau-rect}
 \boxed{
 \tau^{\mathrm H}_{p,\mathcal S}(\xi)
 =\frac{2}{m^2}\left[
 \xi Q_{\mathcal S}-B_{\mathcal S}J_p(\xi)
 +s_*I^m_\xi(\xi)
 \right].}
\end{equation}
\end{theorem}

The large-prime term in \eqref{eq:tau-rect} is deliberately rectangular: it
charges every vulnerable pivot at the largest possible pole grade $s_*$.  A
stronger theorem records the finite $\ell$-adic thickness available above
each pole grade.

For $1\le b\le s_*$ put
\begin{equation}\label{eq:rho-b}
 \rho_b=\sum_{\substack{s\in\mathcal S\\s\ge b}}s(s-b+1).
\end{equation}
Let
\[
 w_m(x)=K_x+\positivepart{m-(K_x+1)x},
 \qquad
 K_x=\floor{\frac{m-1}{x}},
\]
and define the capped window
\begin{equation}\label{eq:I-cap-def}
 I^{m,\rho}_\xi(\xi)=\int_\xi^m\min\{w_m(x),\rho\}\,\dd x.
\end{equation}

\begin{theorem}[Thickness-sensitive denominator polygon]
\label{thm:polygon}
Under the hypotheses of \cref{thm:multiweight}, joint rationality also forces
\begin{equation}\label{eq:poly-holonomy}
 m\left(\LambdaAH(\boldsymbol\phi)
 -\tau^{\mathrm{poly}}_{p,\mathcal S}(\xi)\right)
 \le\Energy(\boldsymbol\phi),
\end{equation}
where
\begin{equation}\label{eq:tau-poly}
 \boxed{
 \tau^{\mathrm{poly}}_{p,\mathcal S}(\xi)
 =\frac{2}{m^2}\left[
 \xi Q_{\mathcal S}-B_{\mathcal S}J_p(\xi)
 +\sum_{b=1}^{s_*}I^{m,\rho_b}_\xi(\xi)
 \right].}
\end{equation}
Moreover
\begin{equation}\label{eq:I-cap-explicit}
 I^{m,\rho}_\xi(\xi)=
 \begin{cases}
 \displaystyle
 \left(m-\frac12\right)H_\rho-\rho\xi,
 &1<\xi<\dfrac{m}{\rho+1},\\[3mm]
 I^m_\xi(\xi),&\xi\ge\dfrac{m}{\rho+1}.
 \end{cases}
\end{equation}
In particular, $\tau^{\mathrm{poly}}_{p,\mathcal S}(\xi)
\le\tau^{\mathrm H}_{p,\mathcal S}(\xi)$.
\end{theorem}

The counting statements are certified from \cref{thm:multiweight}, not from
the refinement.

\begin{theorem}[Explicit simultaneous irrationality bounds]
\label{thm:counts}
Among the $1200$ values
\[
 \zeta_p(3),\zeta_p(5),\ldots,\zeta_p(2401),
\]
the following numbers are irrational:
\begin{center}
\begin{tabular}{@{}crrr@{}}
\toprule
$p$&max. rational&min. irrational&cutoff for $1000$\\
\midrule
$2$&$6$&$1194$&$s\le2013$\\
$3$&$9$&$1191$&$s\le2019$\\
$5$&$15$&$1185$&$s\le2029$\\
$7$&$20$&$1180$&$s\le2039$\\
$13$&$139$&$1061$&$s\le2279$\\
\bottomrule
\end{tabular}
\end{center}
The cutoff in the last column means that at least $1000$ of the positive odd
values with argument at most that cutoff are irrational.  The first four rows
are certified by the branch-and-bound program; the $p=13$ row follows from
the elementary uniform estimate in \cref{prop:p13-finite}.
\end{theorem}

\begin{theorem}[Certified positive-genus finite-packet bounds]
\label{thm:positive-genus-counts}
For the positive-genus prime levels $p\le43$, the following bounds hold among
\[
 \zeta_p(3),\zeta_p(5),\ldots,\zeta_p(2401).
\]
The quantity $\bar d$ is the safe upper bound
$\bar d=2g_p$ for the finite-map degree $d_p$ supplied by
\cref{thm:finite-map-zero}.
\begin{center}
\begin{tabular}{@{}crrrrr@{}}
\toprule
$p$&$g_p$&$\bar d$&max. rational&min. irrational&cutoff for $1000$\\
\midrule
$11$&$1$&$2$&$190$&$1010$&$s\le2379$\\
$17$&$1$&$2$&$226$&$974$&$s\le2459$\\
$19$&$1$&$2$&$237$&$963$&$s\le2483$\\
$23$&$2$&$4$&$377$&$823$&$s\le2823$\\
$29$&$2$&$4$&$414$&$786$&$s\le2921$\\
$31$&$2$&$4$&$425$&$775$&$s\le2951$\\
$37$&$2$&$4$&$457$&$743$&$s\le3039$\\
$41$&$3$&$6$&$595$&$605$&$s\le3445$\\
$43$&$3$&$6$&$606$&$594$&$s\le3483$\\
\bottomrule
\end{tabular}
\end{center}
The last column means that the prefix ending at the displayed odd argument
contains at least $1000$ irrational values.  These bounds use only the exact
all-prime capacity pair, the general-curve arithmetic type, and the elementary
packet majorants of \cref{prop:finite-packet-certificate}; they do not use an
unproved low-degree map or a genus-zero Hasse saving.
\end{theorem}

For every prime $p$, put
\begin{equation}\label{eq:Rp-def}
 R_p(X)=\#\{3\le s\le X:\ s\text{ odd and }\zeta_p(s)\in\Q\}.
\end{equation}

\begin{theorem}[Quantitative genus-zero density one]\label{thm:density-genus-zero}
For $p\in\{2,3,5,7,13\}$, one has
\begin{equation}\label{eq:density-genus-zero}
 \boxed{
 R_p(X)\le
 \left(\frac{e(p-1)}{12\log p}+o(1)\right)(\log X)^2.}
\end{equation}
Consequently
\[
 \#\{3\le s\le X:\ s\text{ odd and }\zeta_p(s)\notin\Q\}
 =\frac X2-O_p((\log X)^2),
\]
and the irrational values have natural density one among the positive odd
arguments.
\end{theorem}

\begin{theorem}[Explicit density-one irrationality at every fixed prime]
\label{thm:density-all-primes}
Let $d_p$ be the degree of the finite map in
\cref{thm:finite-map-zero}; one may take $d_p=1$ when $g_p=0$ and
$d_p\le2g_p$ when $g_p\ge1$.  For every fixed prime $p$,
\begin{equation}\label{eq:density-all-primes}
 \boxed{
 R_p(X)\le
 \left(\frac{e\,d_p(p-1)}{12\log p}+o_p(1)\right)(\log X)^2.}
\end{equation}
In particular, when $g_p\ge1$,
\begin{equation}\label{eq:density-all-primes-genus}
 R_p(X)\le
 \left(\frac{e\,g_p(p-1)}{6\log p}+o_p(1)\right)(\log X)^2.
\end{equation}
Consequently
\[
 \#\{3\le s\le X:\ s\text{ odd and }\zeta_p(s)\notin\Q\}
 =\frac X2-O_p((\log X)^2),
\]
and the irrational values have natural density one among the positive odd
arguments.  At genus zero, $d_p=1$ and the leading coefficient agrees with
\eqref{eq:density-genus-zero}.
\end{theorem}

The induction begins with the following individual theorem, which is also a
special case of the single-weight machinery developed below.

\begin{theorem}[Base individual values]\label{thm:individual}
The following values are irrational:
\[
 \zeta_2(s)\quad(s=3,5,\ldots,29),
 \qquad
 \zeta_3(s)\quad(s=3,5,\ldots,11),
\]
\[
 \zeta_5(3),\quad\zeta_5(5),\quad\zeta_7(3).
\]
\end{theorem}

\begin{warning}[Logical scope]\label{warn:scope}
All asymptotics are for a fixed finite packet $\mathcal S$ while $D\to\infty$.
The constants and exceptional primes may depend on $p$ and $\mathcal S$.
This is sufficient for every counting contradiction: after assuming that a
specified number of rational values exist, one fixes one such finite set and
then runs the asymptotic argument.
\end{warning}

\section{Published inputs and normalization}\label{sec:inputs}

We use the following external results.

\begin{enumerate}[label=\textup{(E\arabic*)},leftmargin=3em]
\item \textbf{Overconvergent Eisenstein families.}
For $s=2n+1$, Calegari's specialization at the character
$\kappa(a)=a^{1-s}$ gives a weight-$(1-s)$ overconvergent eigenform whose
nonconstant coefficients are $\sigma_{-s}^{(p)}(n)$ and whose constant term is
$\zeta_p(s)/2$; see \cite[Theorem~2.3 and Lemma~2.4]{Calegari2005}.

\item \textbf{Finite-slope continuation.}
Buzzard's theorem extends a finite-slope overconvergent form across the
wide-open $W_0(p)$; see \cite[Theorem~5.2]{Buzzard2003}.  The relevant radii
for $p=2,3,5,7,13$ are recalled in \cref{sec:analytic}.

\item \textbf{Hasse invariants.}
We use Katz's geometric Hasse invariant and $q$-expansion principle
\cite{Katz1973}, together with
\[
 E_{\ell-1}\equiv1\pmod\ell,
 \qquad E_{\ell+1}\equiv E_2\pmod\ell.
\]

\item \textbf{De Rham frame torsor and Gauss--Manin connection.}
We use the Hodge-compatible $B$-torsor and connection of
\cite[Propositions~A.14--A.15]{GPRJ2025}.  The descent to one fixed coarse
modular curve and the uniform spreading-out statements used below are proved
internally.

\item \textbf{Arithmetic holonomy.}
We use Bost's slopes inequality \cite{Bost2001}; the capacitary metrics
and model-theoretic Dirichlet problem of \cite{BostChambertLoir2009}; the
theta Hilbert--Samuel, overflow, and direct-image divisor formulas of
\cite{BostCharles2022}; the determinant vector-bundle Hilbert--Samuel theorem
of \cite[Corollary~4.5.2]{ChenMoriwaki2024}; and the zero-estimate,
source-degree, evaluation-height, and prime-window formulas in
\cite{CDT2024,CDT2025}.
Regularity of the modular stack at the cusps is taken from
\cite{Cesnavicius2017}.  The genus-zero fixed-bundle bridge and all
intrinsic-curve passages needed at arbitrary prime level are proved
internally in \cref{sec:bost,sec:curve-zero,sec:qjet,sec:curve-bost}.
\end{enumerate}

Put
\begin{equation}\label{eq:f}
 f_p(\tau)=\left(\frac{\Delta(p\tau)}{\Delta(\tau)}\right)^{1/(p-1)}.
\end{equation}
For $p\in\{2,3,5,7,13\}$ this is a Hauptmodul at the cusp
$P=\infty$.  For $p=13$, Maier's normalization \cite[(3.2) and Table~4]{Maier2009}
gives the canonical Hauptmodul
$t_{13}=13[13]^2/[1]^2=13f_{13}$ and its degree-$14$ map to the
$j$-line.  In all five
cases
\begin{equation}\label{eq:f-product}
 f_p(q)=q\prod_{n\ge1}
 \left(1+q^n+\cdots+q^{(p-1)n}\right)^{24/(p-1)}
 \in q+q^2\Z[[q]].
\end{equation}
Its inverse satisfies
\begin{equation}\label{eq:q-inverse}
 q(f_p)\in f_p+f_p^2\Z[[f_p]].
\end{equation}
Throughout the genus-zero sections we write $f=f_p$.  In the
arbitrary-prime sections no Hauptmodul is used; all local arithmetic is
performed in the width-one Tate parameter $q$.  In either notation,
\[
 D=q\frac{\dd}{\dd q}.
\]

For admissible adelic continuation templates
$\boldsymbol\phi=(\phi_v)_v$, centered at $P$, define
\begin{equation}\label{eq:Lambda}
 \LambdaAH(\boldsymbol\phi)=\sum_v\log\abs{\phi_v'(0)}_v
\end{equation}
and
\begin{equation}\label{eq:Energy}
 \Energy(\boldsymbol\phi)=
 \sum_{v\mid\infty}\iint_{\T_v^2}
 \log\abs{\phi_v(z)-\phi_v(w)}_v\,\dd m(z)\dd m(w)
 +\sum_{v\nmid\infty}\sup_{z\in\T_v}\log\abs{\phi_v(z)}_v.
\end{equation}
These are normalized as in \cite[Section~7.3]{CDT2024} and
\cite[Section~2]{CDT2025}.

\section{Single-weight Eisenstein--Eichler blocks}\label{sec:blocks}

Fix an odd $s\ge3$.  Define
\begin{equation}\label{eq:J-def}
 J_{p,s}(q)=\sum_{n\ge1}\sigma_{-s}^{(p)}(n)q^n,
 \qquad
 \sigma_{-s}^{(p)}(n)=\sum_{\substack{d\mid n\\p\nmid d}}d^{-s},
\end{equation}
and, for a scalar $\eta$,
\[
 K_{p,s,\eta}=J_{p,s}+\eta.
\]
At $\eta=\zeta_p(s)/2$, input (E1) identifies this with the distinguished
overconvergent form of weight $1-s$.

\begin{lemma}[$U_p$ eigenvalue]\label{lem:Up}
For every $n\ge1$,
\[
 \sigma_{-s}^{(p)}(pn)=\sigma_{-s}^{(p)}(n).
\]
Consequently
\[
 U_pK_{p,s,\zeta_p(s)/2}=K_{p,s,\zeta_p(s)/2}.
\]
\end{lemma}

\begin{proof}
The divisors of $pn$ prime to $p$ are exactly the divisors of $n$ prime to
$p$.  The usual $q$-expansion definition of $U_p$ gives the assertion,
including the constant term.
\end{proof}

For $1\le e\le s$ put
\begin{equation}\label{eq:R-def}
 R_{s,e}=D^{s-e}J_{p,s}.
\end{equation}
When no confusion is possible we write $R_e$.

\begin{lemma}[Raw integration depth]\label{lem:raw-depth}
For every $n\ge1$,
\begin{equation}\label{eq:raw-depth}
 n^e[q^n]R_{s,e}
 =\sum_{\substack{d\mid n\\p\nmid d}}(n/d)^s\in\Z.
\end{equation}
If $L_N=\lcmop(1,\ldots,N)$, then
\[
 L_N^eR_{s,e}(q(f))\pmod{f^{N+1}}
 \in\Z[[f]]/(f^{N+1}).
\]
\end{lemma}

\begin{proof}
The coefficient of $q^n$ is $n^{s-e}\sigma_{-s}^{(p)}(n)$, which gives
\eqref{eq:raw-depth}.  Since $n\mid L_N$ for $n\le N$, the $q$-integrality
follows.  The integral tangent-to-identity substitution \eqref{eq:q-inverse}
uses only coefficients of index at most the target index and preserves the
cumulative cutoff.
\end{proof}

\begin{lemma}[Exact divided Frobenius]\label{lem:Dwork}
Let $\ell\ne p$ be prime.  For $1\le e\le s$,
\begin{equation}\label{eq:Dwork-q}
 R_{s,e}(q)-\ell^{-e}R_{s,e}(q^\ell)\in\Z_\ell[[q]].
\end{equation}
After $q=q(f)$,
\begin{equation}\label{eq:Dwork-f}
 R_{s,e}(q(f))-\ell^{-e}R_{s,e}(q(\varphi_\ell(f)))
 \in\Z_\ell[[f]],
\end{equation}
where
\[
 \varphi_\ell(f)=f(q(f)^\ell).
\]
Moreover, outside a fixed exceptional set,
\begin{equation}\label{eq:Frob-expansion}
 \varphi_\ell(f)=f^\ell+\ell\Delta_\ell(f),
 \qquad
 \Delta_\ell(f)\in f^{\ell+1}\Z_\ell[[f]].
\end{equation}
\end{lemma}

\begin{proof}
If $\ell\mid n$, the $q^n$ coefficient of the left side of
\eqref{eq:Dwork-q} is
\[
 n^{s-e}\left(\sigma_{-s}^{(p)}(n)
 -\ell^{-s}\sigma_{-s}^{(p)}(n/\ell)\right)
 =n^{s-e}\sum_{\substack{d\mid n\\p\nmid d\\\ell\nmid d}}d^{-s},
\]
which is $\ell$-integral; the case $\ell\nmid n$ is immediate.  Substitution
gives \eqref{eq:Dwork-f}.  Finally, in characteristic $\ell$,
$f(q)^\ell=f(q^\ell)$.  Substituting $q=q(f)$ gives
$\varphi_\ell(f)\equiv f^\ell\pmod\ell$; the two sides have the same order
and leading coefficient, proving \eqref{eq:Frob-expansion}.
\end{proof}

Let $\calH$ be the rank-two de Rham bundle of the universal elliptic curve and
put
\[
 V_s=\Sym^{s-1}\calH.
\]
In a local oper frame $v_0,\ldots,v_{s-1}$ with
\[
 \nabla_\partial v_j=(s-1-j)v_{j+1},
 \qquad v_s=0,
\]
define
\begin{equation}\label{eq:BGG-def}
 \operatorname{spl}_s(F)=
 \sum_{j=0}^{s-1}(-1)^j\frac{(s-1)!}{(s-1-j)!}
 \partial^{s-1-j}F\,v_j.
\end{equation}

\begin{lemma}[Canonical BGG lift]\label{lem:BGG}
The expression \eqref{eq:BGG-def} glues to an algebraic differential operator
\[
 \operatorname{spl}_s:\omega^{1-s}\longrightarrow V_s
\]
and satisfies
\begin{equation}\label{eq:BGG}
 \nabla_{V_s}\operatorname{spl}_s(F)=\iota_s(\operatorname{Bol}_sF).
\end{equation}
On the Tate disc, $\operatorname{Bol}_sF=D^sF$.
\end{lemma}

\begin{proof}
Write $c_j=(-1)^j(s-1)!/(s-1-j)!$.  In
$\nabla_\partial\operatorname{spl}_s(F)$, the coefficient of
$\partial^{s-j}Fv_j$ for $j\ge1$ is
$c_j+(s-j)c_{j-1}=0$.  Only the top-oper term remains.  The formula is
characterized by this property and its bottom quotient, hence glues.
The Tate normalization is the standard dual-BGG/Bol normalization; compare
\cite[Proposition~2.4.3]{RodriguezCamargo2023} and
\cite[Proposition~2.4.1]{Urban2013}.
\end{proof}

Put
\begin{equation}\label{eq:Eevil}
 E_{p,s+1}^{\mathrm{evil}}=D^sJ_{p,s}.
\end{equation}
It is the algebraic $p$-stabilized Eisenstein form of weight $s+1$ and is
independent of $\eta$.

\begin{proposition}[Split Eisenstein--Eichler connection]
\label{prop:block-connection}
On
\[
 \calE_{p,s}=\calO\oplus V_s
\]
define
\begin{equation}\label{eq:block-connection}
 \nabla_{\calE}(a,v)=
 \left(\dd a,\nabla_{V_s}v-a\,\iota_s(E_{p,s+1}^{\mathrm{evil}})\right).
\end{equation}
This is an algebraic regular-singular connection, and
\begin{equation}\label{eq:block-leaf}
 S_{p,s,\eta}=(1,\operatorname{spl}_s(K_{p,s,\eta}))
\end{equation}
is horizontal on the Tate disc.
\end{proposition}

\begin{proof}
The connection is algebraic because $E_{p,s+1}^{\mathrm{evil}}$ is algebraic;
integrability is automatic on a curve.  Equation \eqref{eq:BGG} gives
horizontality.
\end{proof}

\begin{proposition}[Single-weight functional rank]\label{prop:block-rank}
The complex monodromy orbit of $S_{p,s,\eta}$ spans a space of dimension
$s+1$.
\end{proposition}

\begin{proof}
The period-polynomial cocycle is nonzero.  Otherwise one could subtract a
polynomial of degree at most $s-1$ from the Eichler primitive and obtain a
moderate-growth modular form of weight $1-s<0$, hence zero, contradicting
$D^sK=E_{p,s+1}^{\mathrm{evil}}\ne0$.  The span of the cocycle values is
$\Gamma_0(p)$-stable.  Since $\Gamma_0(p)$ is Zariski dense in
$\mathrm{SL}_2$ and $\Sym^{s-1}(\C^2)$ is irreducible, that span is the full
$s$-dimensional representation.  Adjoining the scalar coordinate gives rank
$s+1$.
\end{proof}

\section{The simultaneous packet and its zero estimate}\label{sec:simultaneous}

Assume from now on that $\mathcal S$ is fixed and that
\begin{equation}\label{eq:joint-rationality}
 \eta_s=\frac{\zeta_p(s)}2\in\Q
 \qquad(s\in\mathcal S).
\end{equation}
Define
\begin{equation}\label{eq:sim-connection}
 \calE_{p,\mathcal S}=\calO\oplus\bigoplus_{s\in\mathcal S}V_s
\end{equation}
with
\begin{equation}\label{eq:sim-connection-formula}
 \nabla_{\mathcal S}(a,(v_s)_s)=
 \left(\dd a,
 \left(\nabla_{V_s}v_s-a\,\iota_s(E_{p,s+1}^{\mathrm{evil}})\right)_s
 \right).
\end{equation}
Then
\begin{equation}\label{eq:sim-leaf}
 S_{p,\mathcal S}=
 \left(1,(\operatorname{spl}_s(K_{p,s,\eta_s}))_{s\in\mathcal S}\right)
\end{equation}
is horizontal.

\begin{theorem}[Joint functional rank]\label{thm:joint-rank}
The complex monodromy orbit of $S_{p,\mathcal S}$ spans the full fibre of
$\calE_{p,\mathcal S}$.  Its dimension is
\[
 m=1+\sum_{s\in\mathcal S}s.
\]
\end{theorem}

\begin{proof}
Fix a basepoint and put
\[
 V_{\mathcal S}=\bigoplus_{s\in\mathcal S}V_s.
\]
Let $W\subset V_{\mathcal S}$ be the span of all differences
\[
 \rho_{\mathcal S}(\gamma)S_{p,\mathcal S}-S_{p,\mathcal S}.
\]
The scalar coordinate of each difference is zero.  For
$\alpha,\gamma\in\Gamma_0(p)$,
\[
 \rho_V(\alpha)
 \left(\rho_{\mathcal S}(\gamma)S-S\right)
 =\left(\rho_{\mathcal S}(\alpha\gamma)S-S\right)
 -\left(\rho_{\mathcal S}(\alpha)S-S\right),
\]
so $W$ is $\Gamma_0(p)$-stable.  Projection to the $s$-block is the span of
the orbit differences of the single-weight leaf; by
\cref{prop:block-rank}, it is all of $V_s$.

The stabilizer of a linear subspace is Zariski closed, so density of
$\Gamma_0(p)$ makes $W$ an $\mathrm{SL}_2$-subrepresentation.  The
representations $V_s=\Sym^{s-1}(\C^2)$ are pairwise nonisomorphic irreducibles
because the weights in $\mathcal S$ are distinct.  Their direct sum is
multiplicity-free.  A subrepresentation surjecting onto every summand must be
the entire direct sum.  Hence $W=V_{\mathcal S}$.  One branch of $S$ has
scalar coordinate $1$, so it is independent of $W$ and completes the span.
\end{proof}

\begin{corollary}[Joint functional independence]\label{cor:joint-independence}
In any rational algebraic frame, the $m$ coordinate functions of
$S_{p,\mathcal S}$ are linearly independent over the function field
$\C(X_0(p))$.  For $p\in\mathcal P_0$ this is $\C(f_p)$, and hence the
coordinates are also independent over $\Q(f_p)$.
\end{corollary}

\begin{proof}
A rational covector over $\C(X_0(p))$ annihilating one branch is
single-valued, hence annihilates all analytic continuations.  At a generic
basepoint those continuations span the full fibre by \cref{thm:joint-rank};
the covector is zero.
\end{proof}

Choose a fixed divisor $B$, disjoint from $P$, clearing all fixed poles, and
put
\begin{equation}\label{eq:source}
 M_D=H^0\left(X_0(p),\calE_{p,\mathcal S}^\vee(DP+B)\right).
\end{equation}

\begin{corollary}[Genus-zero joint Bost pivot set]\label{cor:joint-zero}
For every $\varepsilon>0$,
\begin{equation}\label{eq:joint-zero}
 \#V_D=mD+O_{p,\mathcal S}(1),
 \qquad
 V_D\subset[0,(m+\varepsilon)D+O_{p,\mathcal S,\varepsilon}(1)].
\end{equation}
\end{corollary}

\begin{proof}
Since $X_0(p)\simeq\mathbf P^1$ for the five primes in $\mathcal P_0$,
Riemann--Roch gives $\dim M_D=mD+O(1)$.  A fixed rational frame and fixed
normalizer embed $M_D$ into
\[
 (\Q[f]_{\le D+C})^m.
\]
The coordinate functions are holonomic and independent over $\C(f)$ by
\cref{cor:joint-independence}.  The Chudnovsky--Osgood estimate in
\cite[Theorem~3.2.13]{CDT2024} bounds the order of a nonzero evaluation by
$(m+\varepsilon)D+O(1)$.  Evaluation is injective, and every nonzero graded
piece of the one-variable vanishing filtration has dimension one, so the
number of jumps equals the source rank.
\end{proof}

\section{Global sources and the blockwise Hasse saving}\label{sec:hasse}

The split bundle yields
\begin{equation}\label{eq:source-split}
 M_D=M_D^0\oplus\bigoplus_{s\in\mathcal S}M_{D,s}^+,
\end{equation}
where
\[
 M_D^0=H^0(X_0(p),\calO(DP+B)),
 \qquad
 M_{D,s}^+=H^0(X_0(p),V_s^\vee(DP+B)).
\]
Riemann--Roch gives
\begin{equation}\label{eq:block-source-rank}
 \rank M_{D,s}^+=sD+O_{p,s}(1).
\end{equation}
After enlarging one fixed exceptional set, this decomposition and the relevant
filtrations are saturated over every good $\Z_\ell$.

We recall the blockwise Hasse mechanism in a form that can be summed over
$\mathcal S$.  Write $s=2r+1$ and
\[
 H=D\log f.
\]
On a fixed algebraic Hodge splitting, let $X$ be Urban's unipotent coordinate
and put $x=X/H$.

\begin{proposition}[Single unipotent action]\label{prop:single-action}
For every $D\gg1$ and every $\lambda\in M_{D,s}^+$, the multiplier of the
normalized undifferentiated row $H^rK_{p,s,\eta_s}$ in the genuine scalar
contraction has the form
\begin{equation}\label{eq:q0}
 q_{0,s,\lambda}(f,x)=\sum_{j=0}^{s-1}Q_{j,s,\lambda}(f)x^j.
\end{equation}
There are a fixed polynomial $g_s(f)$, with $g_s(0)$ a unit, and a constant
$C_s$ such that
\begin{equation}\label{eq:q0-degree}
 g_sQ_{j,s,\lambda}\in\Z[\Sigma^{-1}][f],
 \qquad
 \deg(g_sQ_{j,s,\lambda})\le D+C_s.
\end{equation}
Only one degree-$(s-1)$ unipotent matrix occurs.
\end{proposition}

\begin{proof}
In the monomial basis of $\Sym^{s-1}\calH$, the matrix of
$\left(\begin{smallmatrix}1&x\\0&1\end{smallmatrix}\right)$ and its dual have
entries in $\Z[x]$ of degree at most $s-1$.  Keep the source covector in the
fixed algebraic dual frame and transport the BGG section once.  Equivariance
of contraction moves that one unipotent matrix to either factor; it does not
create a second action.  The undifferentiated $K$ occurs only in the bottom
extremal BGG component.  Multiplication by $f^D$ clears the variable pole at
$P$, while one fixed polynomial clears all remaining poles and the fixed
normalization by $H^{-r}$.
\end{proof}

On the chosen splitting,
\begin{equation}\label{eq:x-sigma}
 x_\sigma=-\frac{E_2}{12H}+r_\sigma(f),
 \qquad r_\sigma(f)\in\Q(f).
\end{equation}

\begin{lemma}[Hasse algebraization]\label{lem:Hasse-degree}
For every sufficiently good auxiliary prime $\ell$, the reduction of
$x_\sigma$ lies in $\F_\ell(f)$.  If
\[
 \bar x_\sigma=\frac{a_\ell(f)}{b_\ell(f)}
\]
is in lowest terms and $b_\ell(0)\ne0$, then
\begin{equation}\label{eq:Hasse-degree}
 \delta_\ell:=\max\{\deg a_\ell,\deg b_\ell\}
 \le c_p\ell+C_s.
\end{equation}
\end{lemma}

\begin{proof}
Modulo $\ell$, represent $E_2$ by $E_{\ell+1}/E_{\ell-1}$.  Equation
\eqref{eq:x-sigma} becomes a weight-zero modular function.  The valence formula
on $X_0(p)$ bounds the zero and pole degree by
$[\mathrm{SL}_2(\Z):\Gamma_0(p)](\ell+1)/12=(p+1)(\ell+1)/12$; the fixed
rational correction and elliptic-point conventions change only $O_{p,s}(1)$.
Regularity at the Tate cusp gives $b_\ell(0)\ne0$.
\end{proof}

\begin{theorem}[Blockwise genuine-source Hasse kernel]
\label{thm:block-Hasse}
There are constants $D_0,C_s$, independent of $D$ and of every good $\ell$,
such that the reduction of $M_{D,s}^+$ contains a subspace
$\overline K_{D,\ell,s}$ of dimension at least
\begin{equation}\label{eq:block-Hasse-rank}
 (s-1)\max\{0,D-\ceil{c_p\ell}-C_s\}
\end{equation}
for which
\begin{equation}\label{eq:q0-zero}
 q_{0,s,\bar\lambda}
 \left(f,\frac{a_\ell}{b_\ell}\right)=0
 \quad\text{in }\F_\ell[[f]].
\end{equation}
It lifts to a saturated direct summand $K_{D,\ell,s}\subset M_{D,s,\ell}^+$
with
\begin{equation}\label{eq:q0-divisible}
 q_{0,s,\lambda}(f,x_\sigma(f))\in\ell\Z_\ell[[f]].
\end{equation}
The construction is independent of $\eta_s$.
\end{theorem}

\begin{proof}
After multiplying by $g_s b_\ell^{s-1}$, substitution of
$x=a_\ell/b_\ell$ in \eqref{eq:q0} defines an $\F_\ell$-linear map
\[
 \overline M_{D,s}^+
 \longrightarrow
 \F_\ell[f]_{\le D+(s-1)\delta_\ell+O_s(1)}.
\]
The source has dimension $sD+O_s(1)$, while the target has dimension at most
$D+(s-1)\delta_\ell+O_s(1)$.  Rank--nullity and
\cref{lem:Hasse-degree} give \eqref{eq:block-Hasse-rank}.  Since
$g_s(0)b_\ell(0)$ is a unit, vanishing after clearing denominators implies
\eqref{eq:q0-zero}.  Lift a basis and extend it to an integral basis to obtain
a saturated summand.  The connection and multiplier are independent of the
constant $\eta_s$.
\end{proof}

The genuine contraction separates the unique depth-$s$ row from the remaining
rows:
\begin{equation}\label{eq:deepest-decomp}
 f^D\langle\lambda,S_{p,s,\eta_s}\rangle
 =G_{s,\lambda}+q_{0,s,\lambda}H^rJ_{p,s}
 +\eta_sq_{0,s,\lambda}H^r
 +\sum_{e=1}^{s-1}\widetilde Q_{s,e,\lambda}R_{s,e}.
\end{equation}
Every remaining nonalgebraic row has depth at most $s-1$.

Fix $1<\xi<m$ and put
\[
 N^\sharp=\floor{\xi D}+C_{\mathrm{sh}},
 \qquad a_{\ell,D}=v_\ell(L_{N^\sharp}).
\]
For each good small prime and each $s\in\mathcal S$, choose a complement
$K^c_{D,\ell,s}$ and define the block lattice
\begin{equation}\label{eq:block-hybrid-lattice}
 M_{D,\ell,s}^{\mathrm H}
 =\ell^{sa_{\ell,D}-1}K_{D,\ell,s}
 \oplus\ell^{sa_{\ell,D}}K^c_{D,\ell,s}.
\end{equation}
Together with the unscaled scalar block and the direct sum over $s$, this is
the multiweight small-prime lattice.

\begin{proposition}[Multiweight small-prime cost]\label{prop:small-cost}
The lattice \eqref{eq:block-hybrid-lattice} has integral evaluation through
$f^{\floor{\xi D}}$, and
\begin{equation}\label{eq:small-cost}
 \sum_{\ell\le N^\sharp}
 \log[M_{D,\ell}:M^{\mathrm H}_{D,\ell}]
 =\left(\xi Q_{\mathcal S}-B_{\mathcal S}J_p(\xi)\right)D^2+o(D^2).
\end{equation}
The residual small-prime evaluation heights are $o(D^2)$.
\end{proposition}

\begin{proof}
On $K_{D,\ell,s}$, \eqref{eq:q0-divisible} restores the one missing power of
$\ell$ on the unique depth-$s$ row.  Since
$sa_{\ell,D}-1\ge ea_{\ell,D}$ for $e\le s-1$, every other row is integral as
well.  The block index is
\[
 s^2Da_{\ell,D}
 -(s-1)\max\{0,D-\ceil{c_p\ell}-C_s\}
 +O_s(a_{\ell,D}).
\]
Sum over $s\in\mathcal S$.  The prime number theorem and partial summation give
\[
 \sum_{\ell\le\xi D}a_{\ell,D}\log\ell=\xi D+o(D)
\]
and
\[
 \sum_{\ell\le\xi D}
 \max\{0,D-c_p\ell\}\log\ell
 =D^2J_p(\xi)+o(D^2).
\]
The finite exceptional set and proper-prime-power residuals contribute only
$O_{p,\mathcal S}(D\log D)=o(D^2)$.
\end{proof}

\section{Multiweight Cartier strictness at large primes}\label{sec:Cartier}

On one fixed flat torsor chart, the normalized simultaneous contraction has
the form
\begin{equation}\label{eq:multi-torsor}
 F_\lambda(f)=H_\lambda(f,u)+
 \sum_{s\in\mathcal S}\sum_{e=1}^s
 Q_{s,e,\lambda}(f,u)R_{s,e}(f),
\end{equation}
where every multiplier has $f$-support in $[-C,D+C]$ and fixed fibre degree.
Inserting \cref{lem:Dwork} gives
\begin{equation}\label{eq:multi-Dwork}
 F_\lambda=H_\lambda^\flat+
 \sum_{s\in\mathcal S}\sum_{e=1}^s
 \ell^{-e}Q_{s,e,\lambda}(f,u)
 R_{s,e}(\varphi_\ell(f)),
\end{equation}
with integral $H_\lambda^\flat$.

\begin{lemma}[Coefficient-ring Taylor no-backflow]\label{lem:no-backflow}
Let $A$ be an $\F_\ell$-algebra and
\[
 \Gamma(f,z)=\sum_{j=0}^df^jP_j(z),
 \qquad d<\ell,
 \qquad P_j(z)\in A[[z]].
\]
If $\Gamma(f,f^\ell)\in f^nA[[f]]$ and
$\Delta(f)\in f^{\ell+1}A[[f]]$, then for every $t\ge0$,
\begin{equation}\label{eq:no-backflow}
 \Delta(f)^t
 \left.\partial_z^t\Gamma(f,z)\right|_{z=f^\ell}
 \in f^{n+t}A[[f]].
\end{equation}
\end{lemma}

\begin{proof}
Write $P_j(z)=\sum_rp_{j,r}z^r$.  Since $0\le j<\ell$, the exponent
$j+\ell r$ uniquely determines $(j,r)$, so no cancellation between different
pairs is possible, even when $A$ has zero divisors.  A nonzero coefficient
therefore satisfies $j+\ell r\ge n$.  After $t$ derivatives and multiplication
by $\Delta^t$, its order is at least
\[
 j+\ell(r-t)+t(\ell+1)=j+\ell r+t\ge n+t.
\]
\end{proof}

Let
\[
 F_\lambda(f)=c_nf^n+O(f^{n+1}),\qquad c_n\ne0,
\]
be a Bost pivot.  Since $n=O_{\mathcal S}(D)$ by
\cref{cor:joint-zero} and $\ell>\xi D$, terms of $z$-degree at least $\ell$ in
\eqref{eq:multi-Dwork} lie beyond the pivot range.  Truncate every raw row
below $z^\ell$ and shift by $f^C$.  Put
\begin{equation}\label{eq:Bse}
 B_{s,e}(f,u,z)=f^CQ_{s,e,\lambda}(f,u)R_{s,e,<\ell}(z).
\end{equation}
Then
\[
 \deg_fB_{s,e}<\ell,
 \qquad \deg_zB_{s,e}<\ell,
 \qquad B_{s,e}(f,u,0)=0.
\]
Choose digits
\begin{equation}\label{eq:digits}
 B_{s,e}\equiv\sum_{\nu=0}^{e-1}\ell^\nu B_{s,e,\nu}
 \pmod{\ell^e}.
\end{equation}
For $1\le b\le s_*$ define
\begin{equation}\label{eq:Gamma-b}
 \Gamma_b(f,z)=
 \sum_{s\in\mathcal S}\sum_{e=b}^s
 \overline B_{s,e,e-b}(f,u,z).
\end{equation}
If $\Pi_b$ denotes the pole-grade-$b$ symbol of the shifted truncated
expression, Taylor expansion gives
\begin{equation}\label{eq:pole-symbol}
 \boxed{
 \Pi_b(f)=\sum_{t=0}^{s_*-b}
 \frac{\overline\Delta_\ell(f)^t}{t!}
 \left.\partial_z^t\Gamma_{b+t}(f,z)\right|_{z=f^\ell}.}
\end{equation}
All factorials are units after excluding the finite set $\ell\le s_*$.

\begin{proposition}[Multiweight torsor Cartier strictness]
\label{prop:multi-Cartier}
If $v_\ell(c_n)<0$ and $a=-v_\ell(c_n)$, then
\[
 1\le a\le s_*
\]
and there exist
\[
 -C\le j\le D+C,
 \qquad r\ge1,
 \qquad n=j+\ell r.
\]
Consequently
\begin{equation}\label{eq:rect-local}
 \delta_\ell(\psi_D^{(n)})
 \le s_*v_\ell
 \left(\lcmop\{\max(n-D-C,1),\ldots,n+C\}\right)\log\ell.
\end{equation}
\end{proposition}

\begin{proof}
After truncation, the only negative powers of $\ell$ are the explicit
$\ell^{-e}$ with $e\le s_*$, so $a\le s_*$.  Every coefficient below the
pivot order vanishes.  Taking pole-grade symbols gives
$\Pi_b\in f^{n+C}A[[f]]$.  Descending induction on $b$, using
\cref{lem:no-backflow} in \eqref{eq:pole-symbol}, proves
\[
 \Gamma_b(f,f^\ell)\in f^{n+C}A[[f]],
 \qquad
 \Pi_b\equiv\Gamma_b(f,f^\ell)\pmod{f^{n+C+1}}.
\]
At the actual pole grade $a$, the coefficient of $f^{n+C}$ in $\Pi_a$ is the
nonzero reduction of $\ell^ac_n$.  Hence the same coefficient in
$\Gamma_a(f,f^\ell)$ is nonzero.  Some contributing monomial has
\[
 n+C=j'+\ell r,
 \qquad0\le j'\le D+2C,
 \qquad r\ge1.
\]
Set $j=j'-C$.  Then $\ell r=n-j$ lies in the displayed interval.  Since
$a\le s_*$, \eqref{eq:rect-local} follows.
\end{proof}

\begin{proposition}[Rectangular large-prime cost]\label{prop:rect-large}
The total positive denominator penalty at primes $\ell>N^\sharp$ satisfies
\begin{equation}\label{eq:rect-large}
 \mathscr A_D^{\mathrm{large}}
 \le s_*I^m_\xi(\xi)D^2+o(D^2).
\end{equation}
\end{proposition}

\begin{proof}
Apply \cref{prop:multi-Cartier} and the CDT packing sum for the
$mD+O(1)$ distinct pivots contained in an interval of normalized length
$m+o(1)$.  Fixed shifts by $C$ and $C_{\mathrm{sh}}$ contribute
$O(D\log D)=o(D^2)$.
\end{proof}

\subsection{The thickness-sensitive Fitting refinement}

Let $L_{D,\ell}$ be the integral simultaneous source lattice at a good large
prime and
\[
 L_{D,\ell}^{(n)}=
 L_{D,\ell}\cap\{\lambda:\ord_f\langle\lambda,S\rangle\ge n\}
\]
its saturated vanishing filtration.  At a pivot $n$, write
\[
 a_{\ell,n}=\max\{0,-v_\ell(c_n)\}.
\]
For $1\le b\le s_*$ define
\begin{equation}\label{eq:Qb}
 \mathcal Q_{b,\ell}=
 \bigoplus_{\substack{s\in\mathcal S\\s\ge b}}
 L^+_{D,s,\ell}/\ell^{s-b+1}L^+_{D,s,\ell}.
\end{equation}
Its length is
\[
 \operatorname{length}_{\Z_\ell}\mathcal Q_{b,\ell}
 =\rho_bD+O_{p,\mathcal S}(1),
\]
with $\rho_b$ as in \eqref{eq:rho-b}.

\begin{proposition}[Pole-grade Fitting bound]\label{prop:Fitting}
For every good large prime $\ell$ and $1\le b\le s_*$,
\begin{equation}\label{eq:Fitting}
 \#\{n\in V_D:a_{\ell,n}\ge b\}
 \le\rho_bD+O_{p,\mathcal S}(1).
\end{equation}
\end{proposition}

\begin{proof}
Let $U_{b,n}$ be the image of $L_{D,\ell}^{(n)}$ in
$\mathcal Q_{b,\ell}$.  These form a descending chain.  Suppose
$a_{\ell,n}\ge b$, and choose a primitive lift $\lambda_n$ of the rank-one
pivot quotient.  If its image in $\mathcal Q_{b,\ell}$ belonged to
$U_{b,n+1}$, there would be $\mu\in L_{D,\ell}^{(n+1)}$ such that
\[
 (\lambda_n-\mu)_s\in\ell^{s-b+1}L^+_{D,s,\ell}
 \qquad(s\ge b).
\]
The $s$-block has maximum pole depth $s$, so these components now have depth
at most $b-1$; blocks with $s<b$ already have depth at most $b-1$, and the
scalar block has depth zero.  Since $\mu$ vanishes through the pivot order,
$\lambda_n-\mu$ has the same leading coefficient $c_n$, forcing
$v_\ell(c_n)\ge-(b-1)$, a contradiction.  Therefore
$U_{b,n+1}\subsetneq U_{b,n}$.  Each grade-$\ge b$ pivot consumes at least one
unit of the total length of $\mathcal Q_{b,\ell}$, proving
\eqref{eq:Fitting}.
\end{proof}

\begin{remark}[Why rank alone is insufficient]
One depth-$s$ source direction may successively consume several powers of
$\ell$.  This is why the correct cap is the length
$s-b+1$ per direction, not merely the rank of the deep source.  The quotient
\eqref{eq:Qb} records exactly that thickness.
\end{remark}

\begin{proposition}[Capped prime-window sum]\label{prop:capped-window}
For every integer $\rho\ge1$, the function \eqref{eq:I-cap-def} satisfies
\eqref{eq:I-cap-explicit}.  Moreover the pole-grade Fitting bounds imply
\begin{equation}\label{eq:poly-large}
 \mathscr A_D^{\mathrm{large}}
 \le\left(\sum_{b=1}^{s_*}I^{m,\rho_b}_\xi(\xi)\right)D^2+o(D^2).
\end{equation}
\end{proposition}

\begin{proof}
At a prime $\ell=xD$, the usual degree-$D$ multiplier window contains at most
$w_m(x)D+o(D)$ pivots.  By \cref{prop:Fitting}, at most
$\rho_bD+o(D)$ of them can have pole grade at least $b$.  Hence the grade-$b$
contribution is bounded by
$\min\{w_m(x),\rho_b\}D+o(D)$.  Prime-number-theorem summation yields
\eqref{eq:poly-large}.

The function $w_m$ is decreasing and equals $\rho$ at
$x=m/(\rho+1)$.  If $\xi$ is below that point, the cap contributes
$\rho(m/(\rho+1)-\xi)$ before the ordinary tail begins.  Evaluating
$I^m_\xi$ at $m/(\rho+1)$ gives \eqref{eq:I-cap-explicit}.  If the cutoff is
already beyond the crossing, the cap is inactive.
\end{proof}

\section{Bost--slopes assembly}\label{sec:bost}

We record the fixed-bundle bridge needed for literal application of the scalar
CDT estimates.

\begin{lemma}[Adelic fixed-bundle comparison]\label{lem:fixed-bundle}
Let $\calV$ be a fixed rank-$m$ vector bundle on $\mathbf P^1_\Q$ with fixed
adelic models and metrics, and put
\[
 N_D=H^0(\mathbf P^1,\calV(DP+B)).
\]
Suppose scalar contraction with a fixed horizontal leaf is injective and its
pivot set satisfies a linear zero estimate.  Then fixed changes of frame,
normalizer, bundle splitting, local model, or scalar row degree change the
source degree and summed evaluation heights by $o(D^2)$.  In particular,
\begin{equation}\label{eq:fixed-degree}
 \widehatdeg N_D=\frac m2\Energy D^2+o(D^2),
\end{equation}
and
\begin{equation}\label{eq:fixed-heights}
 \sum_{n\in V_D}\sum_vh_v(\psi_D^{(n)})
 \le\left(-\frac{m^2}{2}\LambdaAH+m\Energy\right)D^2
 +\mathscr A_D+o(D^2).
\end{equation}
\end{lemma}

\begin{proof}
By Birkhoff--Grothendieck, after a fixed twist
$\calV\simeq\bigoplus_i\calO(a_i)$.  The original and transported adelic norms
are uniformly equivalent; the determinant distortion is $O(D)$.  The scalar
arithmetic Hilbert--Samuel formula gives
$\frac12\Energy(D+a_i)^2$ on each summand, proving
\eqref{eq:fixed-degree}.  Fixed analytic matrices and normalizers contribute
$O(\log D)$ per pivot, hence $O(D\log D)=o(D^2)$ in total.  The scalar CDT
height estimate contributes $-n\LambdaAH+D\Energy$ at each pivot.  Since the
$mD+O(1)$ pivot orders are distinct and nonnegative,
\[
 \sum_{n\in V_D}n\ge\frac12m^2D^2+o(D^2),
\]
which proves \eqref{eq:fixed-heights}.
\end{proof}

\begin{proof}[Proof of \cref{thm:multiweight}]
By \cref{prop:small-cost,prop:rect-large,lem:fixed-bundle}, the multiweight
Hasse lattice satisfies
\[
 \widehatdeg M_D^{\mathrm H}
 =\left(\frac m2\Energy-
 \xi Q_{\mathcal S}+B_{\mathcal S}J_p(\xi)\right)D^2+o(D^2)
\]
and
\[
 \sum_{n\in V_D}\sum_vh_v(\psi_D^{(n)})
 \le\left(-\frac{m^2}{2}\LambdaAH+m\Energy
 +s_*I^m_\xi(\xi)\right)D^2+o(D^2).
\]
Bost's slopes inequality gives the first expression no larger than the
second.  Rearranging yields \eqref{eq:multiweight-holonomy} with
\eqref{eq:tau-rect}.
\end{proof}

\begin{proof}[Proof of \cref{thm:polygon}]
Replace \cref{prop:rect-large} by \cref{prop:capped-window}.  The small-prime
lattice and every other Bost term are unchanged.  The same rearrangement gives
\eqref{eq:tau-poly}.
\end{proof}

\section{Common analytic continuation and collision energy}\label{sec:analytic}

We first close the extra continuation interface required at $p=13$.

\begin{lemma}[The $13$-adic supersingular annulus]\label{lem:p13-annulus}
Put $t=13f_{13}$.  Then the unique supersingular annulus of $X_0(13)$ is
\[
 0<v_{13}(t)<1.
\]
The wide-open $W_0(13)$ containing the Tate cusp is therefore
\begin{equation}\label{eq:W013}
 W_0(13)=\{v_{13}(t)>0\}=\{|f_{13}|_{13}<13\}.
\end{equation}
\end{lemma}

\begin{proof}
By \cite[(3.2) and Table~4]{Maier2009}, Maier's canonical Hauptmodul is
$t=t_{13}=13f_{13}$ and the degree-$14$ map to the $j$-line is
\begin{equation}\label{eq:j13}
 j=\frac{(t^2+5t+13)
 (t^4+7t^3+20t^2+19t+1)^3}{t}.
\end{equation}
Deuring's count \cite{Deuring1941} gives a single supersingular
$j$-invariant in characteristic $13$.  It is $j=5$: the curve
\[
 E:y^2=x^3+x+4
\]
has $j(E)=5$ in $\F_{13}$, and the exact quadratic-character sum
\[
 \sum_{x\in\F_{13}}\left(\frac{x^3+x+4}{13}\right)=0
\]
gives $\#E(\F_{13})=14$.  Its Frobenius trace is therefore zero, so $E$ is
supersingular.

Write $v=v_{13}(t)$.  If $v<0$, the highest-degree terms in the two factors
of the numerator of \eqref{eq:j13} dominate, and
$v_{13}(j)=13v<0$.  If $0<v<1$, the terms $5t$ and $1$ dominate, so
$v_{13}(j)=0$ and
\[
 j\equiv5\pmod{13}.
\]
If $v>1$, the terms $13$ and $1$ dominate, and
$v_{13}(j)=1-v<0$.  Thus the open tube with supersingular reduction is
$0<v<1$; the loci $v=0$ and $v=1$ are its two boundary ends and are not part
of the strict wide-open.  The cusp $\infty$ satisfies $v(t)>1$; adjoining the
supersingular annulus to its ordinary component gives $v(t)>0$, which is
\eqref{eq:W013} because $t=13f_{13}$.
\end{proof}

\begin{theorem}[Common continuation]\label{thm:continuation}
For $p\in\{2,3,5,7,13\}$ and every finite packet $\mathcal S$, the
simultaneous horizontal leaf extends over the same domains as its individual
blocks.  At the distinguished $p$-adic place it extends over
\[
 \abs{f_p}_p<p^{12/(p-1)}.
\]
At every finite place $\ell\ne p$, under joint rationality it is meromorphic
on $\abs{f_p}_\ell<1$.  At the complex place it is meromorphic on every
modular template used below.
\end{theorem}

\begin{proof}
Each block has $U_p$-eigenvalue $1$ by \cref{lem:Up}.  Inputs (E1)--(E2)
extend it across the Buzzard wide-open $W_0(p)$.  For $p=2,3$ the explicit radii are the calculations of
\cite[Lemma~3.1 and the proof of Theorem~3.5]{Calegari2005}.  For
$p=5,7$ the identification of that wide-open with
$|f_p|_p<p^{12/(p-1)}$ is the one used in the underlying single-weight
argument.  For $p=13$ it is \cref{lem:p13-annulus}.  The BGG operator applies
only finitely many Gauss--Manin derivatives and does not shrink the base
domain.  A finite direct sum has the intersection of the blockwise domains,
which here is the same common domain.
\end{proof}

For $0<Y<1/p$, put
\[
 \phi_{p,Y}(z)=f_p(e^{-2\pi Y}z)
\]
and
\begin{equation}\label{eq:Cp}
 C_p(Y)=4\pi
 \sum_{\substack{c\ge1\\p\mid c\\cY<1}}
 \frac{\varphi(c)}{c^2}
 \left(\arccos(cY)-cY\sqrt{1-c^2Y^2}\right).
\end{equation}
The expression $C_p(Y)$ itself makes sense for every prime;
\cref{thm:allprime-capacity-pair} below identifies it intrinsically as the
Archimedean modular collision term at arbitrary prime level.

\begin{proposition}[Modular Jensen formula]\label{prop:Jensen}
For every $p\in\{2,3,5,7,13\}$ one has
\begin{equation}\label{eq:Jensen}
 \iint_{\T^2}\log\abs{\phi_{p,Y}(z)-\phi_{p,Y}(w)}\,\dd m(z)\dd m(w)
 =-2\pi Y+C_p(Y).
\end{equation}
Consequently
\begin{equation}\label{eq:analytic-data}
 \Lambda_p(Y)=\frac{12\log p}{p-1}-2\pi Y,
 \qquad
 \Energy_p(Y)=\Lambda_p(Y)+C_p(Y).
\end{equation}
\end{proposition}

\begin{proof}
For all five primes $f_p$ is a Hauptmodul.  The nonidentity collisions are
therefore parametrized by primitive bottom rows $(c,d)$ with $c>0$ and
$p\mid c$.  For
$\gamma=\left(\begin{smallmatrix}a&b\\c&d\end{smallmatrix}\right)$,
\[
 \Im(\gamma(x+iY))=\frac{Y}{(cx+d)^2+c^2Y^2}.
\]
Jensen contributes $2\pi(\Im(\gamma\tau)-Y)_+$.  Unfolding the sum over
primitive $d$ and evaluating the elementary integral gives the summand in
\eqref{eq:Cp}.  The leading $q$ contributes $-2\pi Y$, while the eta-product
factors have logarithmic mean zero.  Adding the distinguished $p$-adic radius
from \cref{thm:continuation} gives \eqref{eq:analytic-data}.
\end{proof}

Combining \cref{thm:multiweight,prop:Jensen}, joint rationality forces
\begin{equation}\label{eq:margin}
 (m-1)\Lambda_p(Y)-m\tau^{\mathrm H}_{p,\mathcal S}(\xi)-C_p(Y)\le0.
\end{equation}
The branch-and-bound certificates prove strict positivity of the left side
for every forbidden aggregate state at $p=2,3,5,7$; the $p=13$ finite example
below uses a separate elementary bound.

\section{Explicit individual and counting applications}\label{sec:applications}

\subsection{The base individual theorem}

For completeness, \cref{tab:individual} records exact rational witnesses for
\cref{thm:individual}.  The margins are rigorous outward interval lower
bounds from the single-weight certificate.

\begingroup\small
\begin{longtable}{@{}rrrrrl@{}}
\caption{Single-weight witnesses used to start the counting induction.}
\label{tab:individual}\\
\toprule
$p$&$s$&$\xi$&$\tau_{p,s}(\xi)$&$Y$&certified lower margin\\
\midrule
\endfirsthead
\toprule
$p$&$s$&$\xi$&$\tau_{p,s}(\xi)$&$Y$&certified lower margin\\
\midrule
\endhead
2&3&$22/19$&$1321/608$&$16/79$&$9.745786633966$\\
2&5&$29/26$&$18445/5616$&$6/49$&$13.454727090366$\\
2&7&$110/101$&$105585/25856$&$1/11$&$15.805524331097$\\
2&9&$89/83$&$10906993/2324000$&$3/43$&$17.190005995811$\\
2&11&$262/247$&$464385991/89631360$&$3/52$&$17.790824360622$\\
2&13&$181/172$&$144989487/25958240$&$2/41$&$17.737379764768$\\
2&15&$478/457$&$16666007411/2810615808$&$3/71$&$17.124944575766$\\
2&17&$305/293$&$12175769323/1954839744$&$3/80$&$16.023242852853$\\
2&19&$758/731$&$11693630797/1801184000$&$1/30$&$14.484310655378$\\
2&21&$461/446$&$536638388131/79764338368$&$1/33$&$12.553794492186$\\
2&23&$1102/1069$&$22108077629089/3185347156992$&$1/36$&$10.264641765760$\\
2&25&$649/631$&$21299386717285/2985462031552$&$1/39$&$7.646580292916$\\
2&27&$1510/1471$&$456079886220493/62372115303680$&$1/42$&$4.724811329943$\\
2&29&$869/848$&$14321093408042087/1915475381820000$&$1/46$&$1.522247623229$\\
3&3&$33/29$&$2029/928$&$4/27$&$6.200294149415$\\
3&5&$87/79$&$112477/34128$&$7/73$&$6.941905556192$\\
3&7&$55/51$&$53435/13056$&$1/15$&$6.163732078651$\\
3&9&$267/251$&$33037241/7028000$&$1/19$&$4.332047506141$\\
3&11&$393/373$&$702154669/135354240$&$3/70$&$1.679644254174$\\
5&3&$11/10$&$177/80$&$7/71$&$2.515479614080$\\
5&5&$29/27$&$12899/3888$&$1/16$&$0.131799356827$\\
7&3&$33/31$&$2219/992$&$1/14$&$0.436413057785$\\
\bottomrule
\end{longtable}
\endgroup

\subsection{The finite branch-and-bound certificate}

For each $p$, list the rational arguments in increasing order as
$r_{p,1}<r_{p,2}<\cdots$.  The individual theorem gives the first lower bound.
At the $k$th stage, fix a proposed upper bound $x$ for $r_{p,k}$.  For each
possible sum $T$ of the previous $k-1$ entries, the margin depends on those
entries only through
\[
 m=1+T+x,
 \qquad Q=\sum_{i<k}r_{p,i}^2+x^2,
 \qquad B=m-1-k,
 \qquad s_*=x.
\]
For fixed $T$ and $x$, the arithmetic type is increasing in $Q$.  The following
majorization lemma therefore reduces the search to one extremal configuration
per aggregate state.

\begin{lemma}[Reverse-greedy square-sum maximizer]\label{lem:convex-search}
Let $a_1<\cdots<a_n<x$ be odd integers with prescribed lower bounds and
prescribed total sum.  Among all such tuples, $\sum a_i^2$ is maximized by
allocating the available excess successively to $a_n,a_{n-1},\ldots$, subject
to the gaps $a_i\le a_{i+1}-2$ and $a_n\le x-2$.
\end{lemma}

\begin{proof}
If $i<j$, $a_i$ can be decreased by $2$, and $a_j$ increased by $2$ without
violating the constraints, then the square sum changes by
\[
 (a_i-2)^2+(a_j+2)^2-a_i^2-a_j^2=4(a_j-a_i)+8>0.
\]
Repeated exchanges force all available excess as far to the right as the
constraints allow, which is exactly the reverse-greedy configuration.
\end{proof}

The certificate evaluates \eqref{eq:margin} for every aggregate state with a
fixed rational pair $(\xi,Y)$ at each stage.  All algebraic quantities are
exact.  The values $\pi$ and $\log p$, square roots, inverse trigonometric
functions, and collision sums are enclosed by outward-rounded fixed-point
intervals.  No binary floating point is used.

The resulting rational-index lower bounds are:
\begin{align*}
 p=2:\quad&31,105,261,553,1069,1937,>2401;\\
 p=3:\quad&13,39,87,169,303,513,831,1301,1983,>2401;\\
 p=5:\quad&7,15,29,51,81,127,189,275,391,547,751,1019,1365,1809,2373,>2401;\\
 p=7:\quad&5,9,15,25,39,57,81,111,153,205,271,355,459,589,749,945,\\
 &\hspace{17mm}1185,1477,1829,2255,>2401.
\end{align*}
Here the $j$th displayed finite number is a certified lower bound for
$r_{p,j}$.

\begin{table}[ht]
\centering
\caption{Explicit counting examples; the first four rows are interval-certified.}
\label{tab:counts}
\begin{tabular}{@{}crrrrr@{}}
\toprule
$p$&aggregate states&smallest stage margin&rational cap&irrational count&$1000$ cutoff\\
\midrule
2&$2{,}591{,}454$&$>2.507569047577$&6&1194&2013\\
3&$4{,}136{,}947$&$>2.005014256508$&9&1191&2019\\
5&$6{,}889{,}305$&$>0.687480581811$&15&1185&2029\\
7&$9{,}493{,}807$&$>0.388266442986$&20&1180&2039\\
13&--&$>0.008$&139&1061&2279\\
\bottomrule
\end{tabular}
\end{table}

\begin{proposition}[A uniform finite bound at $p=13$]\label{prop:p13-finite}
Among
\[
 \zeta_{13}(3),\zeta_{13}(5),\ldots,\zeta_{13}(2401),
\]
at most $139$ values are rational.  Hence at least $1061$ are irrational.  In
the shorter list ending at $s=2279$, at least $1000$ are irrational.
\end{proposition}

\begin{proof}
Suppose that $140$ values are rational and take them as $\mathcal S$.  Then
\[
 m=1+\sum_{s\in\mathcal S}s\ge1+(3+5+\cdots+281)=19881,
 \qquad s_*\le2401,
\]
and $Q_{\mathcal S}\le s_*(m-1)$.  Drop the favorable Hasse subtraction in
\eqref{eq:tau-rect} and let $\xi\downarrow1$.  Since
\[
 I_1^m(1)=\left(m-\frac12\right)H_{m-1}-(m-1),
\]
the terms $s_*(m-1)$ cancel.  Since the right-hand side has this finite
right limit as $\xi\downarrow1$, an admissible $\xi>1$ may be chosen
arbitrarily close to $1$, and
\begin{equation}\label{eq:p13-tau-bound}
 \inf_{\xi>1}\tau^{\mathrm H}_{13,\mathcal S}(\xi)
 \le\frac{2s_*(m-\tfrac12)H_{m-1}}{m^2}.
\end{equation}
The sequence $H_n/(n+1)$ is decreasing for $n\ge2$, since
$H_n>1$.  The classical bound $H_n<\log n+\gamma+1/(2n)$, together with
$\gamma<289/500$ and $\log19880<10$, gives the exact rational estimate
\[
 \inf_{\xi>1}\tau^{\mathrm H}_{13,\mathcal S}(\xi)
 <\frac{4802\cdot10579}{1000\cdot19881}
 <\frac{255523}{100000}.
\]

Choose $Y=1/m$.  Since
$0\le\arccos x-x\sqrt{1-x^2}\le\pi/2$ and
$\varphi(c)\le c$,
\[
 C_{13}(1/m)
 \le\frac{2\pi^2}{13}H_{\lfloor(m-1)/13\rfloor}.
\]
The functions $H_{m-1}/m$ and $(1+\log m)/m$ decrease for the present range;
the first assertion follows from $H_{m-1}>1$, and the second by differentiation.
Using
\[
 \log13>\frac{25649}{10000},\qquad \pi<\frac{22}{7},\qquad
 H_n\le1+\log n,
\]
and again $\log19881<10$, the analytic lower bound is minimized at
$m=19881$ and satisfies the exact rational estimate
\begin{align*}
 \frac{(m-1)\Lambda_{13}(1/m)-C_{13}(1/m)}m
 &\ge \left(1-\frac1{19881}\right)
       \left(\frac{25649}{10000}-\frac{44}{7\cdot19881}\right)\\
 &\qquad-\frac{2(22/7)^2}{13\cdot19881}\cdot11
 >\frac{256361}{100000}.
\end{align*}
The difference between this lower bound and the preceding arithmetic upper
bound is greater than $8/1000$, contradicting \eqref{eq:margin}.  Thus at most $139$ values in the
full list are rational.  The same proof applies to the prefix ending at
$2279$, where $s_*$ is smaller; that prefix contains $1139$ odd arguments, so
at least $1139-139=1000$ are irrational.
\end{proof}

\begin{proof}[Proof of \cref{thm:counts}]
The displayed rational-index bounds show that, up to $2401$, there can be at
most $6,9,15,20$ rational values for $p=2,3,5,7$, respectively.  Subtracting
from $1200$ gives the first four certified rows of \cref{tab:counts}.  For the
$1000$ cutoffs, there are respectively $1006,1009,1014,1019$ positive odd
arguments up to $2013,2019,2029,2039$.  The same rational-index lower bounds
show that at most $6,9,14,19$ of those shorter lists can be rational.  The
$p=13$ row is \cref{prop:p13-finite}.  Hence at least $1000$ are irrational in
each of the five displayed cases.
\end{proof}

\subsection{Certificate files and manifests}\label{sec:certificate-files}

The base individual margins in \cref{tab:individual} are reproduced by
\path{hybrid_rational_interval_certificate.py}; its stored output and manifest
are \path{hybrid_rational_interval_certificate_output.txt} and
\path{hybrid_rational_interval_certificate_sha256.txt}.  The manifest records
\begin{verbatim}
1408c9092d0c34917253c8f520db56853112c58640d15c98d58b76a61a0478f5
  hybrid_rational_interval_certificate.py
9e1d8f74eb198c7e7b0c7420a1e8021d7bad4d2bc8014cc939b6b353111c0da0
  hybrid_rational_interval_certificate_output.txt
\end{verbatim}

The multiweight program
\path{multiweight_all_primes_interval_certificate.py} checks
$23{,}111{,}513$ aggregate states.  Its stored output is
\path{multiweight_all_primes_interval_certificate_output.txt}.  The manifest
\path{multiweight_all_primes_interval_certificate_sha256.txt} records
\begin{verbatim}
4f9cfbfe9691447fb352a000b360029bdef46f4f66dea7fa4a45f414f247c1f0
  multiweight_all_primes_interval_certificate.py
46b3325cd5a9e6e81b7dfa565361d845e1480aa8f7c3251727b81ca4b583e89b
  multiweight_all_primes_interval_certificate_output.txt
\end{verbatim}


The positive-genus exact-capacity bounds in
\cref{thm:positive-genus-counts} are checked by
\path{allprime_capacity_finite_packet_certificate.py}.  The script uses
outward-rounded integer fixed-point intervals only; its stored output is
\path{allprime_capacity_finite_packet_certificate_output.txt}.  The manifest
\path{allprime_capacity_finite_packet_certificate_sha256.txt} records
\begin{verbatim}
52e404ad0ea803526a0f682cb0f9e8049ff34d29f32bbd1c590925d421c1168f
  allprime_capacity_finite_packet_certificate.py
3e2a4104ab71582b5ed1873805c2ccdd9f2d867d54a9669fda1a43be6574f75b
  allprime_capacity_finite_packet_certificate_output.txt
\end{verbatim}

\section{Quantitative density one at genus-zero prime level}\label{sec:density}

Put
\[
 L_p=\frac{12\log p}{p-1}.
\]
We now prove \cref{thm:density-genus-zero}.  Only the rectangular multiweight theorem
is needed; in particular, neither the Hasse subtraction nor the Fitting
polygon is load-bearing.

\begin{proof}[Proof of \cref{thm:density-genus-zero}]
Fix $\beta>1$.  For integers $j\ge1$, let $\mathcal S_j$ be the set of
rational odd arguments in the multiplicative shell
\[
 \beta^j\le s<\beta^{j+1},
\]
and put $k_j=\#\mathcal S_j$.  Empty shells may be ignored.  Set
$A=\beta^j$ and form the simultaneous packet attached to $\mathcal S_j$.
Then
\[
 k_jA\le m-1<k_j\beta A,
 \qquad s_*<\beta A,
 \qquad Q_{\mathcal S_j}\le s_*(m-1)<\beta A m.
\]
Choose $\xi=1+1/\log A$ for $A$ sufficiently large, drop the favorable
Hasse term, and use
\[
 I_\xi^m(\xi)\le I_1^m(1)<mH_{m-1}.
\]
The arithmetic type therefore satisfies
\begin{equation}\label{eq:shell-type}
 \tau^{\mathrm H}_{p,\mathcal S_j}(\xi)
 \le\frac{2\beta}{k_j}\bigl(\xi+H_{m-1}\bigr).
\end{equation}

For the analytic template take $Y=1/m$.  Since
\[
 0\le\arccos x-x\sqrt{1-x^2}\le\frac\pi2
\]
and $\varphi(c)\le c$, the collision formula gives
\begin{equation}\label{eq:collision-log-bound}
 C_p(1/m)
 \le\frac{2\pi^2}{p}H_{\lfloor(m-1)/p\rfloor}
 =O_p(\log m).
\end{equation}
Hence the rationality inequality \eqref{eq:margin} implies
\begin{equation}\label{eq:analytic-shell-lower}
 \tau^{\mathrm H}_{p,\mathcal S_j}(\xi)
 \ge L_p-O_p\left(\frac{\log m}{m}\right).
\end{equation}
Initially $k_j=O_\beta(A)$, so $m=O_\beta(A^2)$ and
$H_{m-1}=O_\beta(\log A)$.  For all sufficiently large $j$,
\eqref{eq:analytic-shell-lower} is at least $L_p/2$; combining it with
\eqref{eq:shell-type} therefore first gives
$k_j=O_{p,\beta}(\log A)$.  Substitution back into the source-rank bound gives
\[
 m=O_{p,\beta}(A\log A),
 \qquad
 H_{m-1}=\log A+O_{p,\beta}(\log\log A).
\]
Consequently
\begin{equation}\label{eq:shell-count}
 k_j\le
 \left(\frac{2\beta}{L_p}+o(1)\right)\log A
 =\left(\frac{2\beta\log\beta}{L_p}+o(1)\right)j.
\end{equation}

If $\beta^J\le X<\beta^{J+1}$, summing
\eqref{eq:shell-count} over $j\le J$ yields
\[
 R_p(X)
 \le\left(\frac{\beta}{L_p\log\beta}+o(1)\right)(\log X)^2.
\]
The function $\beta/\log\beta$ has its minimum $e$ at $\beta=e$.  This proves
\eqref{eq:density-genus-zero}.  There is no hidden uniformity assumption in this
passage: for each fixed shell, the finite packet $\mathcal S_j$ is fixed while
the source degree $D$ tends to infinity; only after the exact multiweight
inequality has been obtained do we let $j$ tend to infinity.  Packet-dependent
exceptional primes and lower-order constants have therefore disappeared before
the shell summation.  Since the number of positive odd integers at most $X$ is
$X/2+O(1)$, the density assertion follows.
\end{proof}

\begin{remark}[Why this does not give all but finitely many]\label{rem:not-cofinite}
The estimate permits an infinite exceptional set with only
$O_p((\log X)^2)$ elements up to $X$.  The multiweight mechanism gains rank
only from values simultaneously assumed rational, so it cannot by itself
exclude an arbitrarily sparse sequence of rational exceptions.  This is the
precise motivation for the anchor-packet problem.
\end{remark}



\section{Arbitrary prime level: finite-map restriction of scalars}
\label{sec:curve-zero}

Fix an arbitrary prime $p$, put $X_p=X_0(p)_{\Q}$, and let $P=\infty$ be
the width-one Tate cusp.  Add the finite elliptic locus and every fixed pole
to a divisor $B$ disjoint from $P$.

\begin{lemma}[Coarse-curve descent at arbitrary prime level]
\label{lem:allprime-descent}
For every finite packet of odd weights, the bundle
\[
 \calO\oplus\bigoplus_{s\in\mathcal S}\Sym^{s-1}\calH,
\]
its simultaneous Eisenstein--Eichler connection, canonical BGG lifts,
horizontal germ, and contraction pairing descend from one fixed neat cover
to $X_p\setminus|B|$.  After enlarging a fixed exceptional set, all these
objects and the source lattices spread over $\Z[\Sigma^{-1}]$, uniformly in
$D$.  The descended horizontal leaf has functional rank
$m=1+\sum_{s\in\mathcal S}s$ and independent coordinates over
$\C(X_p)$.
\end{lemma}

\begin{proof}
The frame torsor and Gauss--Manin structures are functorial on a fixed neat
cover, so their two pullbacks to the fibre square agree and satisfy the cocycle
condition.  Away from the elliptic points the coarse-moduli morphism still has
the residual generic $\{\pm1\}$-inertia; the stack is not literally equal to
its coarse curve.  However, $-1$ acts trivially on
$\Sym^{s-1}\calH$ because $s-1$ is even, and it acts trivially on the scalar
summand, connection, BGG operator, and contraction pairing.  Coherent bundles
with trivial inertia descend along the tame coarse-moduli morphism.  Removing
the finitely many elliptic points eliminates the additional inertia, so these
objects descend to $X_p\setminus|B|$.  The BGG lift is characterized
intrinsically by its bottom quotient and top-oper derivative and is therefore
compatible with descent.  Clearing the finitely many coefficients and fixed
poles gives a model over $\Z[\Sigma^{-1}]$.

The joint monodromy proof of \cref{thm:joint-rank} is genus-independent: each
nonzero Eichler cocycle spans the irreducible module
$\Sym^{s-1}(\C^2)$, and the distinct modules form a multiplicity-free direct
sum.  The functional-independence argument of
\cref{cor:joint-independence} then applies over $\C(X_p)$.
\end{proof}

Let $g_p$ be the genus of $X_p$.  The second cusp $Q=0$ is rational and
distinct from $P$.

\begin{theorem}[Finite-map Shidlovsky zero estimate]
\label{thm:finite-map-zero}
There is a finite map
\[
 \pi_p:X_p\longrightarrow\mathbf P^1_{\Q}
\]
unramified at $P$, of a fixed degree $d_p$, such that the source
\[
 M_D=H^0\!\left(X_p,\calE_{p,\mathcal S}^{\vee}(DP+B)\right)
\]
has
\begin{equation}\label{eq:curve-pivot-count}
 \#V_D=mD+O_{p,\mathcal S}(1)
\end{equation}
and
\begin{equation}\label{eq:curve-pivot-range}
 V_D\subset[0,d_pmD+O_{p,\mathcal S}(1)].
\end{equation}
One may take $d_p=1$ when $g_p=0$ and $d_p\le2g_p$ when $g_p\ge1$.
\end{theorem}

\begin{proof}
Assume $g_p\ge1$ and put $A=(2g_p-1)Q$.  Riemann--Roch gives
$\ell(A+P)=\ell(A)+1$, so there is
$x\in L(A+P)\setminus L(A)$.  It has a simple pole at $P$, defines a finite
map $\pi_p=x$ of degree at most $2g_p$, and is unramified at $P$.  Thus
$t=x^{-1}$ and the Tate parameter $q$ are both uniformizers at $P$ and define
the same order function.

Write $\calF=\calE_{p,\mathcal S}^{\vee}(B)$ and
$\pi_p^*(\infty)=P+E$.  Since $DP\le D\pi_p^*(\infty)$, projection formula
gives an injection
\[
 M_D\hookrightarrow
 H^0\!\left(\mathbf P^1,\pi_{p*}\calF\otimes\calO(D\infty)\right).
\]
The pushforward is a fixed rank-$d_pm$ vector bundle.  After one fixed
Birkhoff--Grothendieck splitting, normalized contractions are polynomial
combinations of degree $D+O(1)$ of $d_pm$ fixed formal functions.

Let $k=\Q(t)$ and $K=\Q(X_p)$, choose a $k$-basis
$y_1,\ldots,y_{d_p}$ of $K$, and extend scalars to $\C$.  If
$f_1,\ldots,f_m$ are the horizontal coordinates, then the functions
$y_jf_i$ are independent over $\C(t)$: a relation would be a
$\C(X_p)$-linear relation among the $f_i$.  The derivation $d/dt$ extends
uniquely to the finite separable extension $K/k$; after scalar extension,
restricting the first-order connection system therefore gives a
derivative-stable $d_pm$-dimensional $\C(t)$-space.  The sharp Shidlovsky
estimate
\cite[Theorem~3.2.8]{CDT2024} bounds every nonzero normalized contraction by
$d_pmD+O(1)$.  Finally, Riemann--Roch gives
$\dim M_D=mD+O(1)$, evaluation is injective by
\cref{lem:allprime-descent}, and the nonzero graded quotients of a one-variable
vanishing filtration are one-dimensional.  This proves both assertions.  The
genus-zero case uses a degree-one coordinate.
\end{proof}

\begin{remark}[Why the covering degree enters only the range]
The map $\pi_p$ multiplies the number of auxiliary holonomic functions by
$d_p$, but it does not multiply the polynomial degree $D$.  The genuine
source retains rank $mD+O(1)$; only the ambient pivot interval grows to
$d_pmD+O(1)$.
\end{remark}

\section{Intrinsic \texorpdfstring{$q$}{q}-jet Cartier strictness}
\label{sec:qjet}

The genus-zero proof uses a Hauptmodul and a Taylor correction for its
coordinate Frobenius.  On $X_p$ we instead remain in the width-one Tate
parameter.  For every auxiliary prime $\ell\ne p$ the exact relation is
\begin{equation}\label{eq:q-Dwork-all}
 R_{s,e}(q)-\ell^{-e}R_{s,e}(q^\ell)\in\Z_\ell[[q]].
\end{equation}

Fix an algebraic local frame of every coefficient bundle at $P$ and use the
canonical local frame $q^{-D}$ of $\calO(DP)$.  If
$h\in H^0(X_p,\calW(DP+B_0))$, its \emph{normalized $q$-germ} is
$q^Dh$ after trivializing $B_0$ at $P$.  One further fixed shift $q^C$ may be
used to remove the bounded negative support coming from the torsor chart.
Every normalized coefficient and every normalized pivot below is understood
in this sense.  This convention makes the order estimates independent of an
unstated choice of pole-clearing coordinate.

\begin{lemma}[Global coefficient bundle and jet injectivity]
\label{lem:q-jet-injective}
For a fixed packet there are a fixed vector bundle $\calW$, a fixed divisor
$B_0$ disjoint from $P$, and a constant $C$ such that every coefficient of a
normalized torsor multiplier is the $q$-germ of a section of
\[
 H^0(X_p,\calW(DP+B_0)).
\]
After excluding a fixed finite set of auxiliary primes, every nonzero reduced
normalized coefficient has order at most $D+C$.  Consequently, for
$\ell>D+C$, reduction followed by the $\ell$-jet map at $P$ is injective.
\end{lemma}

\begin{proof}
The representation, BGG operators, bounded torsor-monomial degree, algebraic
contraction, and frame changes are fixed.  Coefficient extraction is therefore
an algebraic map from the source bundle to a fixed finite-rank coefficient
bundle; poles outside the chosen chart are absorbed in $B_0$.

Choose a fixed saturated injection
$\calW(B_0)\hookrightarrow\calO(A)^{\oplus r}$ with $A$ disjoint from $P$.
For a nonzero scalar component $h\in H^0(X_p,\calO(DP+A))$, write
$h=q^{-D}u(q)$ in the chosen local frame.  If $h$ has a pole at $P$, then
$\ord_qu<D$; if it is regular at $P$, the principal-divisor identity gives
$\ord_Ph\le\deg A$ and hence $\ord_qu\le D+\deg A$.  Thus every nonzero
normalized component has order at most $D+\deg A$.  After spreading out
and excluding finitely many primes, the injection and the relevant section
spaces remain saturated and commute with reduction; the same divisor-degree
argument holds on the smooth geometrically integral special fibre.  A
nonzero section in the kernel of the $\ell$-jet map would then have order at
least $\ell>D+C$, a contradiction.
\end{proof}

\begin{lemma}[Exact $q$-block no-backflow]
\label{lem:q-block}
Let
\[
 \Gamma(q,z)=\sum_{r=0}^{\ell-1}w_r(q)z^r
\]
with coefficients in the reduced coefficient space of
\cref{lem:q-jet-injective}, and suppose $\ell>D+C$.  If
$\Gamma(q,q^\ell)\ne0$, then
\[
 \ord_q\Gamma(q,q^\ell)=j+\ell r_0
\]
for some $0\le j\le D+C$ and $0\le r_0<\ell$.  If the $z$-constant term is
zero, then $r_0\ge1$.
\end{lemma}

\begin{proof}
Let $r_0$ be the least index with $w_{r_0}\ne0$.  Jet injectivity gives
$j=\ord_qw_{r_0}\le D+C<\ell$.  Its first term after substitution occurs
strictly before $\ell(r_0+1)$, whereas every later $z$-block begins at least
at $\ell(r_0+1)$.  No later block can cancel it.  This is valid over an
arbitrary coefficient algebra, including one with zero divisors.
\end{proof}

\begin{proposition}[General-curve $q$-jet Cartier strictness]
\label{prop:q-jet-Cartier}
There are a fixed exceptional set and a constant $C$ such that, for
$D\gg1$ and every good auxiliary prime satisfying
\[
 \ell>D+C,\qquad \ell>s_*,\qquad \ell\ne p,
\]
the following holds.  If a normalized Bost pivot
\[
 F_\lambda(q)=c_nq^n+O(q^{n+1}),\qquad c_n\ne0,
\]
has $v_\ell(c_n)<0$ and $a=-v_\ell(c_n)$, then
\begin{equation}\label{eq:q-pole-grade}
 1\le a\le s_*
\end{equation}
and
\begin{equation}\label{eq:q-window}
 n=j+\ell r,
 \qquad -C\le j\le D+C,
 \qquad r\ge1.
\end{equation}
Consequently
\begin{equation}\label{eq:q-local-height}
 \delta_\ell(\psi_D^{(n)})
 \le s_*v_\ell\!\left(
 \lcmop\{\max(n-D-C,1),\ldots,n+C\}\right)\log\ell.
\end{equation}
\end{proposition}

\begin{proof}
Insert \eqref{eq:q-Dwork-all} into the genuine torsor decomposition.  By
\cref{thm:finite-map-zero}, $n=O(D)$; for $\ell>D+C$ and $D\gg1$ one has
$n+C<\ell^2$.  Terms of raw $z$-degree at least $\ell$ therefore lie beyond
the pivot.  The coefficients below $z^\ell$ are $\ell$-integral, so the only
negative powers are the explicit $\ell^{-e}$ with $e\le s_*$, proving
\eqref{eq:q-pole-grade}.

Digitize each complete truncated product, coefficientwise in the finite free
global coefficient lattice.  For pole grade $b$ let $\Gamma_b(q,z)$ be the
sum of the corresponding digits.  Because the substitution is exactly
$q\mapsto q^\ell$, its associated-graded symbol is simply
\[
 \Pi_b(q)=\Gamma_b(q,q^\ell);
\]
there are no Taylor correction terms.  At the actual pole grade $a$, the
$q^n$-coefficient is the nonzero reduction of $\ell^ac_n$, and every lower
coefficient vanishes.  Thus $\ord_q\Gamma_a(q,q^\ell)=n$.  The raw rows have
zero constant term, so \cref{lem:q-block} gives
\eqref{eq:q-window}.  The positive integer $n-j=\ell r$ lies in the stated
interval, and $a\le s_*$ gives \eqref{eq:q-local-height}.
\end{proof}

For $d,m\ge1$ and $1<\xi<dm$, put
\begin{equation}\label{eq:Gdm}
 G_{d,m}(\xi)=\int_\xi^{dm}\min\left\{1,\frac d x\right\}\,\dd x
 =\begin{cases}
 d-\xi+d\log m,&1<\xi<d,\\[1mm]
 d\log\!\left(\dfrac{dm}{\xi}\right),&d\le\xi<dm.
 \end{cases}
\end{equation}

\begin{proposition}[General-curve large-prime cost]
\label{prop:curve-large}
For the unmodified source at primes $\ell>\xi D$,
\begin{equation}\label{eq:curve-large}
 \mathscr A_D^{\mathrm{large}}
 \le s_*mG_{d_p,m}(\xi)D^2+o(D^2).
\end{equation}
\end{proposition}

\begin{proof}
A prime $\ell=xD$ can affect only pivots lying in degree-$D$ windows ending
at positive multiples of $\ell$.  By \eqref{eq:curve-pivot-range} there are
at most $d_pm/x+O(1)$ such windows, while the total number of pivots is
$mD+O(1)$.  Hence the vulnerable-pivot count is at most
$mD\min\{1,d_p/x\}+o(D)$.  Apply \cref{prop:q-jet-Cartier} and the prime
number theorem, and integrate over $x$.
\end{proof}

\section{Positive capacitary width and normalization-robust slopes}
\label{sec:curve-bost}

The arbitrary-prime density theorem needs only a positive tangent-degree term
whose coefficient is quadratic in the source rank.  It does not require an
exact Bost--Charles self-intersection normalization.  We therefore work
directly with finite-order scalar capacitary estimates and keep every
normalization-dependent term at size $O_p(mD^2)$.

For a place $v$, let $(\Omega_v,o_v)$ be a pointed analytic curve and let
\[
 \iota_v:(\Omega_v,o_v)\longrightarrow (X_p^{\mathrm{an}},P)
\]
be nonconstant and unramified at $o_v$.  Transporting the capacitary norm on
$T_{o_v}\Omega_v$ through $D\iota_v(o_v)$ gives a norm
$\|\cdot\|_{v,\circ}$ on $T_PX_p$.  Put
\begin{equation}\label{eq:robust-width}
 \Lambda_p^\circ=
 \widehatdeg\bigl(T_PX_p,(\|\cdot\|_{v,\circ})_v\bigr).
\end{equation}

\begin{lemma}[Strict capacitary collar]
\label{lem:strict-cap-collar}
Let $K$ be a non-Archimedean local field, let $X/K$ be a smooth projective
curve, and let $P\in X(K)$ extend to a section meeting the special fibre of a
normal projective model at a smooth point of an irreducible component $C_P$.
Let $\Omega_0$ be the integral residue disc at $P$.  Let $\Omega_1$ be a
connected wide-open containing $P$ for which, on an adapted model, the
continuation-side vertical set contains $C_P$.  Then
\[
 \|\cdot\|^{\mathrm{cap}}_{P,\Omega_1}
 <\|\cdot\|^{\mathrm{cap}}_{P,\Omega_0}
\]
on every nonzero vector of $T_PX$.
\end{lemma}

\begin{proof}
Let $U_1=X^{\mathrm{an}}\setminus\Omega_1$.  On an adapted normal model, write
the Bost--Chambert-Loir extension of the effective divisor $P$ as
\[
 \mathscr P_{\Omega_1}=\mathscr P+V,
 \qquad V=\sum_C c_C C.
\]
The divisor $V$ is the unique solution of the graph-Dirichlet problem in
\cite[Proposition~5.1]{BostChambertLoir2009}.  More is proved in the proof of
\cite[Proposition~5.6]{BostChambertLoir2009}: if the effective divisor meets
every connected component of the continuation domain, then
\[
 c_C=0\quad\text{on the excluded side},
 \qquad c_C>0\quad\text{on the continuation side}.
\]
Here $\Omega_1$ is connected and contains $P$, while $C_P$ lies on its
continuation side.  Hence
\[
 c_{C_P}>0.
\]

The strict norm comparison is local at the intersection of the section with
$C_P$.  After clearing the denominator of $c_{C_P}$, choose local parameters
$q,\varpi$ with $\mathscr P=(q)$ and $C_P=(\varpi)$.  Relative to the integral
model, a local model frame of
$\mathcal O(\mathscr P+c_{C_P}C_P)$ is
$\varpi^{-c_{C_P}}q^{-1}$.  Under adjunction, $q^{-1}$ maps to
$\partial/\partial q$, and therefore
\begin{equation}\label{eq:strict-cap-factor}
 \left\|\frac{\partial}{\partial q}\right\|^{\mathrm{cap}}_{P,\Omega_1}
 =|\varpi|^{c_{C_P}}
  \left\|\frac{\partial}{\partial q}\right\|^{\mathrm{int}}_P
 <\left\|\frac{\partial}{\partial q}\right\|^{\mathrm{int}}_P.
\end{equation}
The same calculation after a finite extension, followed by the usual
normalization of local degrees, gives the identical inequality when the
coefficients are merely rational.  Finally,
\cite[Example~5.11]{BostChambertLoir2009} identifies the integral tangent norm
with the capacitary norm of the bare residue disc $\Omega_0$.  This proves the
claim.
\end{proof}

\begin{proposition}[Positive common capacitary width]
\label{prop:positive-width}
For every fixed prime $p$, there is an adelic family
$\boldsymbol\Omega_p^\circ=((\Omega_v,o_v,\iota_v))_v$ such that:
\begin{enumerate}[label=\textup{(\roman*)},leftmargin=2.7em]
\item every finite simultaneous Eisenstein--Eichler leaf continues along
$\iota_v$ for every place $v$;
\item at all but finitely many finite places, $\Omega_v$ is the integral Tate
residue disc and $\iota_v$ is its inclusion;
\item the tangent degree in \eqref{eq:robust-width} is strictly positive:
\[
 \Lambda_p^\circ>0;
\]
\item at the distinguished $p$-adic place there is a larger continuation
wide-open $\Omega_p^+$ such that the affinoid closure of $\Omega_p^\circ$ is
compactly contained in $\Omega_p^+$.  Moreover one may choose a rational
function $u_p$ with polar divisor $a_pP$, $a_p\ge1$, for which
\begin{equation}\label{eq:nested-lemniscate}
 \Omega_p^\circ=\{|u_p|_p>1\},
 \qquad
 \{|u_p|_p\ge1\}\Subset\Omega_p^+\Subset W_0(p).
\end{equation}
At every Archimedean place the image of the closed pointed domain is likewise
contained in a slightly larger continuation domain.
\end{enumerate}
The domains and $\Lambda_p^\circ$ depend only on $p$, not on the packet of
weights.
\end{proposition}

\begin{proof}
Work first on a fixed sufficiently fine tame level.  Buzzard's wide-open
$W_0(p)$ is the continuation region through the Tate component, and
\cite[Theorem~5.2]{Buzzard2003} extends every finite-slope overconvergent
eigenform under consideration across it.  Each specialization used here has
$U_p$-eigenvalue one, and applying the algebraic BGG operator does not shrink
the base domain.

Choose descent-invariant cuts strictly inside all supersingular boundary
annuli and pass to the quotient on the coarse curve.  The resulting connected
wide-open $\Omega_p^\circ\Subset W_0(p)$ contains the tube of the coarse
component through $P$ and a positive annular collar at every boundary end.
Choose a second set of cuts farther out in the same annuli; it gives a larger
connected wide-open $\Omega_p^+$ with
$\overline{\Omega_p^\circ}\Subset\Omega_p^+\Subset W_0(p)$.  The descended
leaf continues on $\Omega_p^+$ for every weight.  By
\cref{lem:strict-cap-collar}, the capacitary tangent metric of
$\Omega_p^\circ$ has a strict gain over the integral Tate metric:
\[
 \lambda_p=-\log
 \frac{\|t\|_{p,\Omega_p^\circ}^{\mathrm{cap}}}
      {\|t\|_{p,\mathrm{int}}}>0.
\]
Applying the collar lemma after descent avoids any cusp-width or ramification
normalization from the auxiliary tame cover.

The complement of $\Omega_p^\circ$ is an affinoid meeting every connected
component.  By \cite[Proposition~5.6]{BostChambertLoir2009}, it is a rational
lemniscate: there is a rational function $u_p$, with polar divisor $a_pP$,
such that its complement is $\{|u_p|_p\le1\}$.  Because the cuts defining
$\Omega_p^\circ$ were chosen strictly inside those defining $\Omega_p^+$,
the affinoid closure $\{|u_p|_p\ge1\}$ is compactly contained in
$\Omega_p^+$.  This proves \eqref{eq:nested-lemniscate}.

At every finite place different from $p$, use the integral Tate residue disc.
At the complex place use the pointed unit disc with
\[
 \iota_\infty(z):q=e^{-2\pi Y}z.
\]
Choose it with a slightly larger concentric disc still contained in the
$q$-coordinate continuation region.  Its transported tangent norm contributes
$-2\pi Y$ to the Arakelov degree.  Choose $Y>0$ with
$2\pi Y<\lambda_p$.  The width-one Tate parameter $q$ is an integral
parameter on the fixed cusp model, so $\partial/\partial q$ generates the
reference tangent lattice at every finite place; together with the unit
complex $q$-disc this reference line has degree zero.  Consequently
\[
 \Lambda_p^\circ\ge\lambda_p-2\pi Y>0.
\]
Every domain and every larger continuation neighborhood depends only on $p$,
and a finite direct sum of blocks continues on these same domains.
\end{proof}

Put
\[
 L=\calO_{X_p}(P),\qquad
 \calF_{\mathcal S}=\calE_{p,\mathcal S}^{\vee}(B),\qquad
 M_D=H^0(X_p,\calF_{\mathcal S}\otimes L^{\otimes D}).
\]
The support of $B$ is disjoint from $P$ and is fixed once the packet is fixed.
On each chosen analytic domain, $\calO(B)$ is treated as a fixed line-bundle
twist: after a fixed power and finitely many residue classes it has a
nonvanishing analytic frame.  This contributes only packet-dependent
lower-order constants.  Fix an adelic model of $L$ and a reference adelic
norm $\|\cdot\|_{0,v}$ on $T_PX_p$; at every finite place take the
reference tangent norm to be the integral Tate norm.  Riemann--Roch
gives
\begin{equation}\label{eq:robust-source-rank}
 r_D:=\rank M_D=mD+O_{p,\mathcal S}(1).
\end{equation}

For the arbitrary-prime argument, choose the reference adelic source lattice
blockwise from fixed algebraic presentations as in the next lemma.  All good
prime coefficient lattices used in the Cartier argument are obtained from the
same presentations; enlarging the packet-dependent exceptional set therefore
introduces no second, incompatible source model.

\begin{lemma}[Fixed-twist source and boundary comparison]
\label{lem:robust-fixed-twist}
There are constants $A_p,B_p\ge0$, depending only on $p$, with the following
property.  For every fixed packet $\mathcal S$ there is a constant
$C_{p,\mathcal S}$ such that
\begin{equation}\label{eq:robust-source-lower}
 \widehatdeg M_D\ge-A_pmD^2+o_{p,\mathcal S}(D^2),
\end{equation}
and, at the distinguished $p$-adic and complex domains,
\begin{equation}\label{eq:robust-boundary-comparison}
 \log\|\langle\lambda,S_{p,\mathcal S}\rangle\|_{\partial\Omega_v}
 \le B_pD+C_{p,\mathcal S}\log(D+2)+\log\|\lambda\|_v.
\end{equation}
The boundary estimate remains valid after replacing any finite local source
lattice by a sublattice.  The resulting loss of source degree is not hidden in
\eqref{eq:robust-source-lower}; it is accounted for by its lattice index.
\end{lemma}

\begin{proof}
Fix once and for all an adelic model of the scalar pair $(X_p,L)$.  After
replacing $L$ by a fixed positive power and treating the finitely many residue
classes as fixed twists, the graded section algebra
\[
 R_p=\bigoplus_{e\ge0}H^0(X_p,L^e)
\]
is generated by finitely many homogeneous scalar sections.  Choose integral
lifts of these generators outside one fixed set of places.  In weighted degree
$e$, the resulting monomials give a surjection from a free adelic module
$P_e$ onto $E_e=H^0(X_p,L^e)$.  Equip $E_e$ with the quotient lattice and
quotient norm.  The number of weighted monomials is polynomial in $e$, every
monomial has Arakelov degree bounded below by $-A_pe-O_p(1)$, and its norm on
each of the two fixed analytic boundaries is at most
$\exp(B_pe+O_p(1))$.  Consequently
\begin{equation}\label{eq:scalar-presentation-bounds}
 \widehat\mu_{\min}(E_e)\ge-A_pe-O_p(\log(e+2)),
 \qquad
 \|E_e\longrightarrow H^0(\partial\Omega_v,L^e)\|_{\mathrm{op}}
 \le\exp(B_pe+O_p(\log(e+2))).
\end{equation}
The first estimate is the finite-generation estimate used in
\cite[Theorem~6.1, (6.8)]{BostChambertLoir2009}; the monomial presentation
also makes the second estimate literal.  Replacing these quotient norms by any
fixed standard model or $L^2$ norm changes only the constants $A_p,B_p$.

Because $L$ is ample, for the fixed bundle $\calF_{\mathcal S}$ there are
fixed integers $a_1,\ldots,a_u$ and a sheaf surjection
\[
 \bigoplus_{j=1}^uL^{-a_j}\longrightarrow\calF_{\mathcal S}.
\]
After tensoring by $L^D$, Serre vanishing for the fixed kernel gives a
surjection
\begin{equation}\label{eq:fixed-presentation}
 \bigoplus_{j=1}^uE_{D-a_j}\longrightarrow M_D
\end{equation}
for $D\gg_{p,\mathcal S}1$.  Spread this one presentation out and use its
quotient lattice and quotient norms.  It is chosen separately on the scalar
block and on each positive block, so it respects the direct-sum source used in
\cref{eq:curve-hybrid-lattice}.  After inverting a finite packet-dependent set,
relative Serre vanishing and cohomology-and-base-change identify this quotient
lattice with the corresponding integral section lattice.  At the remaining
places we retain the presentation lattice itself.  Thus the same lattice is
used in the source-degree, coefficient, and Cartier calculations.

A quotient with its quotient norm has minimum slope no smaller than that of
its source.  Applying \eqref{eq:scalar-presentation-bounds} to
\eqref{eq:fixed-presentation}, the fixed shifts $a_j$, the fixed number $u$ of
summands, and the fixed presentation maps contribute only
$O_{p,\mathcal S}(\log(D+2))$.  Hence
\[
 \widehat\mu_{\min}(M_D)
 \ge-A_pD-O_{p,\mathcal S}(\log(D+2)).
\]
Multiplication by \eqref{eq:robust-source-rank} proves
\eqref{eq:robust-source-lower}.  In particular, the coefficient of $D$ is a
scalar constant depending only on $p$; all packet dependence is lower order
for a packet fixed before $D\to\infty$.

Finally compose \eqref{eq:fixed-presentation} with restriction to a boundary
and contraction with the horizontal leaf.  The scalar restriction norm is
bounded by the second inequality in
\eqref{eq:scalar-presentation-bounds}.  The finitely many presentation maps,
the simultaneous leaf, and a fixed analytic frame for $\calO(B)$ have bounded
norms depending on the packet but not on $D$.  The finite direct sum contributes
at most another packet-dependent constant.  Passing to the quotient therefore
gives \eqref{eq:robust-boundary-comparison}.  If a finite source lattice is
replaced by a sublattice, its gauge norm increases, so the operator-norm upper
bound can only improve.
\end{proof}

Let $V_D$ be the rational pivot set of the vanishing filtration and let
$G_{D,n}$ be the saturated rank-one quotient at a pivot $n$, with its quotient
adelic norms.  Denote its evaluation map by $\psi_D^{(n)}$.  For a finite
place $v\ne p$, equip the graded target with its integral model and set
$\delta_{v,D,n}=\max\{0,\log\|\psi_D^{(n)}\|_v\}$.  Thus
$\delta_{v,D,n}$ is exactly the positive excess local height at $v$.  Put
\[
 \delta_{D,n}=\sum_{\substack{v\nmid\infty\\v\ne p}}
                    \delta_{v,D,n}.
\]
The distinguished $p$-adic operator norm is controlled directly by capacity
and is not counted again in $\delta_{D,n}$.

\begin{lemma}[Finite-order scalar capacitary jet estimate]
\label{lem:robust-cap-jet}
There is a constant $B_p'\ge0$, depending only on $p$, such that, for every
fixed packet, there is a constant $C'_{p,\mathcal S}$ satisfying
\begin{equation}\label{eq:robust-single-jet}
 \widehatdeg G_{D,n}
 \le B_p'D-n\Lambda_p^\circ+\delta_{D,n}
     +C'_{p,\mathcal S}\log(D+2)
\end{equation}
for every pivot $n$ and all sufficiently large $D$.  The same estimate holds
for a finite adelic sublattice of the reference source, with its induced
vanishing filtration and quotient norms.
\end{lemma}

\begin{proof}
The vector bundle is removed before the analytic estimate is applied.
Contraction sends $\lambda$ to an analytic section of
$L^D\otimes\calO(B)$; after using the fixed frame of the second factor it is
a scalar analytic section of $L^D$.  Thus the only vector-bundle input is the
bounded contraction operator in \cref{lem:robust-fixed-twist}.

At the $p$-adic place use the finite-order inequality derived inside the
proof of \cite[Proposition~5.15]{BostChambertLoir2009}, not merely its
limiting canonical-norm conclusion.  Before the limiting step, that proof uses
only a trivialization of a fixed power of $L$ and the maximum principle for
the resulting scalar analytic function.  It does not use that the section
comes from a global scalar linear system, and therefore applies verbatim to
the scalar section obtained after contraction.  On a rational inner
lemniscate a fixed power of $L$ is trivial.  Taking $D$ first in the
corresponding arithmetic progression and treating the remaining residue
classes as fixed twists, the calculation bounds the $n$th divided jet by the
boundary norm times
\[
 \|t\|_{p,\circ}^{\,n}\exp(O_p(D)),
\]
relative to a tangent vector $t$.  At the complex place the same assertion is
ordinary Cauchy on the chosen $q$-disc.  Combining these estimates with
\eqref{eq:robust-boundary-comparison} gives
\begin{equation}\label{eq:robust-local-analytic}
 \sum_{v\in\{p,\infty\}}\log\|\psi_D^{(n)}\|_v
 \le B_pD
 +n\sum_{v\in\{p,\infty\}}
   \log\frac{\|t\|_{v,\circ}}{\|t\|_{0,v}}
 +O_{p,\mathcal S}(\log(D+2)).
\end{equation}
The estimate is valid for the quotient $G_{D,n}$ because its quotient norm is
the infimum of the norms of its lifts.

At every remaining finite place one has
$\log\|\psi_D^{(n)}\|_v\le\delta_{v,D,n}$ by definition.  The coherent
models used to define these quantities are the same models used in the
Cartier calculation below; fixed packet denominators are therefore included
in the finite exceptional-place term there, rather than in an unrecorded
analytic constant.  Since the capacitary and reference tangent norms agree at
all finite places other than $p$, summing all places yields
\begin{equation}\label{eq:robust-height}
 h(\psi_D^{(n)})
 \le B_pD-n\bigl(\Lambda_p^\circ-
                   \widehatdeg(\overline T_0)\bigr)
      +\delta_{D,n}+O_{p,\mathcal S}(\log(D+2)).
\end{equation}

The graded target is
\[
 \overline L_P^{\,D}\otimes
 \overline T_0^{\vee\otimes n}\otimes\overline C_{\mathcal S},
\]
where $\overline C_{\mathcal S}$ is fixed with the packet.  The rank-one degree
identity, in the normalization of
\cite[Theorem~6.1, (6.7)]{BostChambertLoir2009}, gives
\[
 \widehatdeg G_{D,n}
 =D\widehatdeg(\overline L_P)
  -n\widehatdeg(\overline T_0)
  +h(\psi_D^{(n)})+O_{p,\mathcal S}(1).
\]
Substituting \eqref{eq:robust-height} cancels the arbitrary reference tangent
norm and proves \eqref{eq:robust-single-jet}.  This is only a formal change of
reference metric; it is not the target-normal $\Theta$-cancellation proved
in \cref{thm:curve-Bost-exact,app:exact-normalization}.  Passing to a finite sublattice increases its gauge norm and
therefore cannot increase the analytic operator norm, proving the final
assertion.
\end{proof}

Fix $1<\xi<d_pm$ and put
\[
 N=\floor{\xi D},\qquad L_N=\lcmop(1,\ldots,N).
\]
Using the blockwise reference lattices above, define
\begin{equation}\label{eq:curve-hybrid-lattice}
 M_D^{\mathrm{hyb}}=
 M_D^0\oplus\bigoplus_{s\in\mathcal S}L_N^sM_{D,s}^+.
\end{equation}
At the finitely many packet-dependent bad places use the corresponding
integral direct-sum presentation lattice.  We use this presentation lattice
throughout.  If it is compared with, or saturated inside, any fixed geometric
model, the quotient is killed by one packet-dependent integer because the
presentation is fixed; since the rank is $O_{p,\mathcal S}(D)$, the resulting
logarithmic index is $O_{p,\mathcal S}(D)=o(D^2)$.

\begin{proposition}[Compatibility with the hybrid finite-prime lattices]
\label{prop:curve-hybrid-compatibility}
For every fixed packet,
\begin{equation}\label{eq:curve-hybrid-degree}
 \widehatdeg M_D^{\mathrm{hyb}}
 \ge-A_pmD^2-\xi Q_{\mathcal S}D^2
       +o_{p,\mathcal S}(D^2),
\end{equation}
and
\begin{equation}\label{eq:curve-hybrid-delta}
 \sum_{n\in V_D}\delta_{D,n}
 \le s_*mG_{d_p,m}(\xi)D^2+o_{p,\mathcal S}(D^2).
\end{equation}
The scalar capacitary estimate \eqref{eq:robust-single-jet} remains valid for
the quotient norms induced by $M_D^{\mathrm{hyb}}$.
\end{proposition}

\begin{proof}
The positive weight-$s$ block has rank $sD+O_{p,s}(1)$.  Scaling it by
$L_N^s$ therefore has logarithmic index
\[
 s^2D\log L_N+O_{p,s}(\log L_N).
\]
Since $\log L_N=N+o(N)$, summing over the blocks and using
\eqref{eq:robust-source-lower} proves
\eqref{eq:curve-hybrid-degree}.  The fixed bad-place and saturation
adjustments are $o(D^2)$.

We next bound the finite local heights away from $p$.  By
\eqref{eq:curve-pivot-range} and the fixed multiplier shifts, every raw
coefficient index entering a pivot is at most $K_{p,\mathcal S}D$ for a fixed
constant $K_{p,\mathcal S}$.  For a prime $\ell\le N$, $\ell\ne p$, raw depth
and multiplication by $L_N^s$ remove the full $\ell$-denominator unless
$v_\ell(k)>v_\ell(L_N)$ for one of these indices $k$.  Up to the fixed factor
$s_*$, the residual valuation at one pivot is bounded by
\begin{align*}
 &\sum_{\substack{\ell\le N\\\ell\ne p}}
 \bigl(v_\ell(L_{\lfloor K_{p,\mathcal S}D\rfloor})
       -v_\ell(L_N)\bigr)_+\log\ell\\
 &\hspace{18mm}\le
 \sum_{\substack{a\ge2\\N<\ell^a\le K_{p,\mathcal S}D}}\log\ell
 =O_{p,\mathcal S}(D^{1/2}\log D).
\end{align*}
Only proper prime powers occur because $N> D$ for large $D$.  Since
$\#V_D=O_{p,\mathcal S}(D)$, all residual small-prime heights total
$O_{p,\mathcal S}(D^{3/2}\log D)=o(D^2)$.  The distinguished place $p$ is
already controlled by the capacitary Cauchy estimate with the hybrid source
norm and is not included in $\delta_{D,n}$.  At each of the finitely many other bad primes $\ell$, the raw-depth
comparison is uniform in $D$.  Indeed,
\[
 \bigl(v_\ell(L_{\floor{K_{p,\mathcal S}D}})-v_\ell(L_N)\bigr)_+
 =O_{p,\mathcal S,\xi}(1),
\]
because the ratio of the two cutoffs is fixed.  The frame, presentation, and
torsor denominators likewise have bounded $\ell$-valuation.  Hence every
fixed bad place contributes $O_{p,\mathcal S,\xi}(1)$ per pivot, and all
such places together contribute $O_{p,\mathcal S,\xi}(D)=o(D^2)$.  The
separate saturation adjustments in the source degree remain
$O_{p,\mathcal S}(D\log(D+2))=o(D^2)$.

For $\ell>N$, the hybrid lattice equals the reference lattice.  Since
$\xi>1$, for the fixed packet and large $D$ one has $\ell>D+C$ and
$\ell>s_*$.  Thus
\cref{prop:q-jet-Cartier,prop:curve-large} applies; the at most one prime in
the harmless floor interval between $N$ and $\xi D$ contributes $o(D^2)$.
This proves \eqref{eq:curve-hybrid-delta}.

Finally, $M_D^{\mathrm{hyb}}\subset M_D$ at every finite place.  Its gauge
norm is therefore larger, and the analytic operator-norm estimates for the
reference source can only improve.  The vanishing filtration uses the same
rational subspaces and the quotient norms induced by this hybrid lattice, so
the finite-height and capacitary calculations refer to the same rank-one
quotients.  This proves the compatibility assertion.
\end{proof}

\begin{theorem}[Normalization-robust general-curve slopes]
\label{thm:curve-holonomy}
For every fixed prime $p$ there are constants
\[
 \Lambda_p^\circ>0,\qquad C_p>0,
\]
depending only on $p$, such that the following holds.  If a fixed finite packet
$\mathcal S$ is jointly rational, then for every $1<\xi<d_pm$,
\begin{equation}\label{eq:curve-tau}
 \tau^{\mathrm{curve}}_{p,\mathcal S}(\xi)
 =\frac{2}{m^2}\left[
   \xi Q_{\mathcal S}+s_*mG_{d_p,m}(\xi)\right]
\end{equation}
satisfies
\begin{equation}\label{eq:robust-curve-slopes}
 \boxed{
 \tau^{\mathrm{curve}}_{p,\mathcal S}(\xi)
 \ge\Lambda_p^\circ-\frac{C_p}{m}.}
\end{equation}
The constants are independent of the packet.
\end{theorem}

\begin{proof}
Use the vanishing filtration on $M_D^{\mathrm{hyb}}$.  Functional independence
makes scalar contraction injective, so the filtration has exactly
$r_D=mD+O_{p,\mathcal S}(1)$ saturated rank-one quotients.  Their Arakelov
degrees add:
\[
 \widehatdeg M_D^{\mathrm{hyb}}
 =\sum_{n\in V_D}\widehatdeg G_{D,n}.
\]
The pivot orders are distinct nonnegative integers, whence
\begin{equation}\label{eq:robust-pivot-sum}
 \sum_{n\in V_D}n
 \ge\frac{r_D(r_D-1)}2
 =\frac{m^2}{2}D^2+o_{p,\mathcal S}(D^2).
\end{equation}
Summing \eqref{eq:robust-single-jet}, using
\eqref{eq:curve-hybrid-delta}, and observing that the packet-dependent
$O(\log D)$ term at each of $O(D)$ pivots is $o(D^2)$, gives
\[
 \widehatdeg M_D^{\mathrm{hyb}}
 \le B_p'mD^2-\frac{m^2}{2}\Lambda_p^\circ D^2
 +s_*mG_{d_p,m}(\xi)D^2+o_{p,\mathcal S}(D^2).
\]
Compare with \eqref{eq:curve-hybrid-degree}, divide by $D^2$, and let
$D\to\infty$.  Then
\[
 \frac{m^2}{2}\Lambda_p^\circ
 \le(A_p+B_p')m+\xi Q_{\mathcal S}
       +s_*mG_{d_p,m}(\xi).
\]
Multiplication by $2/m^2$ proves
\eqref{eq:robust-curve-slopes} with $C_p=2(A_p+B_p')$.

Every packet-dependent presentation, frame, bad-place, and boundary constant
has disappeared into $o_{p,\mathcal S}(D^2)$ before the packet is varied.
This is why $C_p$ is genuinely packet-independent, despite the absence of any
uniform exceptional-prime set as the packet varies.
\end{proof}

The following theorem is the audited replacement for the former exact-normalization
conjecture.  It proves the required source and jet formulas for a canonical
mixed adelic metric on $X_p$.  Only the associated capacitary tangent line is
realized as a Bost--Charles normal line; no whole-divisor direct-image map is
part of the statement.

\begin{theorem}[Exact mixed capacitary normalization and target-normal cancellation]
\label{thm:curve-Bost-exact}
Fix a prime $p$ and the continuation datum
$\boldsymbol\alpha=\boldsymbol\Omega_p^\circ$ of
\cref{prop:positive-width}.  There are a normal arithmetic model
$\mathscr X_\alpha$ of $X_p$ and a semipositive, nef, and big adelic
Arakelov $\Q$-divisor $\widehat P_\alpha$ extending $P$.  Put
\begin{equation}\label{eq:exact-data}
 \overline L_\alpha=\mathcal O(\widehat P_\alpha),\qquad
 \Theta_\alpha=\widehatdeg P^*\overline L_\alpha,\qquad
 \Energy_\alpha=\widehat P_\alpha^{\,2},\qquad
 \LambdaAH_\alpha=\widehatdeg\overline T_\alpha>0,
\end{equation}
where $\overline T_\alpha$ is the source capacitary tangent line transported
to $T_PX_p$.  These quantities depend only on $p$ and on the fixed
continuation datum.

Let $\mathcal S$ be a fixed jointly rational packet, and give the fixed
coefficient bundle $\calF_{\mathcal S}$ any fixed dominated adelic metric
compatible with its integral models.  Equip
\[
 M_D=H^0(X_p,\calF_{\mathcal S}\otimes\mathcal O(DP))
\]
with the induced adelic supremum norms, replacing each Archimedean norm by
its John ellipsoid norm.  Then
\begin{equation}\label{eq:exact-source}
 \boxed{
 \widehatdeg\overline M_D^\alpha
 =\frac m2\Energy_\alpha D^2+o_{p,\mathcal S}(D^2).}
\end{equation}
The leading term is independent of the fixed metric on
$\calF_{\mathcal S}$.

For a pivot $n$, let $\overline G_{D,n}^\alpha$ be the saturated rank-one
source quotient and let
\[
 \overline J_{D,n}^\alpha
 =P^*\overline L_\alpha^{\,D}
  \otimes\overline T_\alpha^{\vee\otimes n}
  \otimes\overline C_{\mathcal S}
\]
be its natural jet target, where $\overline C_{\mathcal S}$ is fixed.  Denote
by $h_\alpha^{\mathrm{nat}}(\psi_D^{(n)})$ the global operator height with
this target metric.  After scalarizing the target by a nonzero rational
vector, denote the resulting height by
$h_\alpha^{\mathrm{sc}}(\psi_D^{(n)})$; it is independent of the chosen
vector by the product formula.  For every $C>0$, uniformly for all pivots
$0\le n\le CD$,
\begin{align}
 \widehatdeg\overline J_{D,n}^\alpha
 &=D\Theta_\alpha-n\LambdaAH_\alpha+O_{p,\mathcal S}(1),
 \label{eq:exact-target-degree}\\
 h_\alpha^{\mathrm{nat}}(\psi_D^{(n)})
 &\le D(\Energy_\alpha-\Theta_\alpha)+\delta_{D,n}
       +O_{C,p,\mathcal S}(1),
 \label{eq:exact-natural-height}\\
 h_\alpha^{\mathrm{sc}}(\psi_D^{(n)})
 &\le D\Energy_\alpha-n\LambdaAH_\alpha+\delta_{D,n}
       +O_{C,p,\mathcal S}(1),
 \label{eq:exact-scalar-height}\\
 \widehatdeg\overline G_{D,n}^\alpha
 &\le D\Energy_\alpha-n\LambdaAH_\alpha+\delta_{D,n}
       +O_{C,p,\mathcal S}(1).
 \label{eq:exact-graded-degree}
\end{align}
Thus the $D\Theta_\alpha$ contribution of the natural target cancels
exactly with the $-D\Theta_\alpha$ part of the natural evaluation estimate.
In particular,
\begin{equation}\label{eq:exact-degree-sum}
 \sum_{n\in V_D}\widehatdeg\overline G_{D,n}^\alpha
 \le\left(-\frac{m^2}{2}\LambdaAH_\alpha
          +m\Energy_\alpha\right)D^2
    +\mathscr A_D+o_{p,\mathcal S}(D^2).
\end{equation}

For the hybrid source of \eqref{eq:curve-hybrid-lattice}, endowed with these
metrics,
\begin{equation}\label{eq:exact-hybrid-source}
 \widehatdeg\overline M_D^{\alpha,\mathrm{hyb}}
 =\left(\frac m2\Energy_\alpha-\xi Q_{\mathcal S}\right)D^2
  +o_{p,\mathcal S}(D^2).
\end{equation}
Therefore joint rationality forces the exact general-curve inequality
\begin{equation}\label{eq:exact-curve-slopes}
 \boxed{
 \tau^{\mathrm{curve}}_{p,\mathcal S}(\xi)
 \ge\LambdaAH_\alpha-\frac{\Energy_\alpha}{m}.}
\end{equation}
All formulas are invariant under finite base change after normalized Arakelov
degrees are divided by the extension degree.
\end{theorem}

\begin{proof}
The canonical mixed divisor and its nefness and bigness are constructed in
\cref{prop:exact-mixed-divisor}.  The denominator-cleared realization of the
capacitary tangent line is \cref{prop:cap-normal-realization}; the base-change
statement is \cref{lem:exact-base-change}.  The determinant source formula is
\cref{prop:exact-source-formula}, and the uniform jet estimate and exact
$D\Theta_\alpha$ cancellation are
\cref{prop:exact-single-jet}.  Finally,
\cref{prop:exact-hybrid-assembly} gives
\eqref{eq:exact-hybrid-source}, \eqref{eq:exact-degree-sum}, and the slopes
assembly \eqref{eq:exact-curve-slopes}.
\end{proof}

\begin{remark}[Exact versus robust normalization]
The exact theorem strengthens but does not replace
\cref{thm:curve-holonomy}.  The normalization-robust theorem still gives an
independent qualitative proof that $R_p(X)=O_p((\log X)^2)$.  The explicit
leading coefficient in \cref{thm:density-all-primes}, however, uses the exact
modular capacity pair proved below.  The exact theorem removes the coarse
source and boundary constants $A_p$ and $B_p'$ and replaces their contribution
by the intrinsic self-intersection $\Energy_\alpha/m$.
\end{remark}

\begin{remark}[Scope of the realization bridge]
The bridge proved below realizes the full metrized capacitary tangent line as
the normal line of a pseudoconcave formal-analytic surface, after clearing
rational finite-place exponents.  Matching this normal line does not determine
the entire Bost--Chambert-Loir finite vertical divisor and does not by itself
produce a Bost--Charles meromorphic map to $\mathscr X_\alpha$.  Neither
stronger assertion is used.  The exact source and jet formulas are proved
directly on $X_p$ for the canonical mixed finite-capacitary/Archimedean-direct-
image metric.
\end{remark}


\subsection{The exact modular capacity pair at every prime}
\label{sec:allprime-capacity}

For $0<Y<1/p$, let the Archimedean pointed disc be
\begin{equation}\label{eq:allprime-qdisc}
 \iota_Y:\mathbb D\longrightarrow X_p(\C),
 \qquad q=e^{-2\pi Y}z,
\end{equation}
and retain the integral Tate residue disc at every finite place $\ell\ne p$.
At the distinguished place $p$, let $W_0(p)$ be Buzzard's maximal wide-open
through the Tate component.

\begin{lemma}[Rational skeleton exhaustion and exact finite capacity]
\label{lem:maximal-capacity-exhaustion}
For every prime $p$ there is an increasing sequence of connected wide-opens in
$X_{p,\C_p}^{\mathrm{an}}$,
\[
 \Omega_{p,1}\Subset\Omega_{p,2}\Subset\cdots\Subset W_0(p),
 \qquad \bigcup_{\nu\ge1}\Omega_{p,\nu}=W_0(p),
\]
such that every $\Omega_{p,\nu}$ descends to a finite extension of $\Q_p$,
has affinoid complement, and is a rational wide-open: after that finite
extension there is a rational function $u_{p,\nu}$, with polar divisor a
positive multiple of $P$, for which
\[
 \Omega_{p,\nu}=\{|u_{p,\nu}|_p>1\}.
\]
If
\begin{equation}\label{eq:inner-capacity-gain}
 \lambda_{p,\nu}=
 -\log\frac{\|\partial_q\|^{\mathrm{cap}}_{p,\Omega_{p,\nu}}}
              {\|\partial_q\|_{p,\mathrm{int}}},
\end{equation}
then
\begin{equation}\label{eq:inner-capacity-limit}
 \lambda_{p,\nu}\longrightarrow
 L_p=\frac{12\log p}{p-1}.
\end{equation}
Moreover $\Omega_{p,\nu+1}$ is a larger continuation neighborhood of the
affinoid closure of $\Omega_{p,\nu}$, so every sufficiently large member,
together with any fixed Archimedean $q$-disc of positive total tangent degree,
is admissible for \cref{thm:curve-Bost-exact}.

For $p\ge5$, after a finite unramified extension splitting the supersingular
residue points, put
\[
 \vartheta_\nu=\frac{\nu}{\nu+1}.
\]
The domain $\Omega_{p,\nu}$ may be chosen by retaining, in the branch indexed
by a geometric supersingular elliptic curve $E/\overline{\F}_p$, the initial
skeleton segment of length $\vartheta_\nu i(E)$, where
\[
 i(E)=\frac{\#\operatorname{Aut}(E)}2.
\]
Its effective resistance and tangent gain are then
\begin{equation}\label{eq:inner-resistance}
 R_{p,\nu}
 =\left(\sum_{[E]\ \mathrm{ss}}
       \frac1{\vartheta_\nu i(E)}\right)^{-1}
 =\vartheta_\nu\frac{12}{p-1},
 \qquad
 \lambda_{p,\nu}=R_{p,\nu}\log p.
\end{equation}
For $p=2,3$, one may instead take
\begin{equation}\label{eq:wild-inner-discs}
 \Omega_{p,\nu}
 =\{|f_p|_p<p^{\vartheta_\nu 12/(p-1)}\},
 \qquad
 \lambda_{p,\nu}
 =\vartheta_\nu\frac{12\log p}{p-1}.
\end{equation}
All normalized local and global Arakelov degrees are independent of the finite
extensions used to define these cuts.
\end{lemma}

\begin{proof}
Assume first that $p\ge5$.  Pass to a finite unramified extension over which
all supersingular residue points split.  By
\cite[\S2, citing Deligne--Rapoport VI.6.16]{ColemanMcMurdy2006}, the regular
semistable skeleton between the two principal ordinary components consists of
parallel paths indexed by the supersingular elliptic curves and having lengths
$i(E)$.  After a further finite ramified extension whose ramification index is
divisible by $\nu+1$, the point at distance $\vartheta_\nu i(E)$ on every
path is a vertex of a semistable subdivision.  Retain the component through
$P$ and the initial path segments up to those vertices.  The tube of this
connected subgraph is $\Omega_{p,\nu}$; its complement is a connected affinoid
neighborhood of the opposite ordinary component.  The inclusions are compact,
and the retained lengths tend to the full branch lengths, so these domains
increase to $W_0(p)$.  This agrees with the affinoid exhaustion of the maximal
wide-open on a fine tame cover in \cite[Proposition~4.1]{Buzzard2003}.

The domain $\Omega_{p,\nu}$ is connected and contains $P$, while its complement
$U_{p,\nu}$ is affinoid.  Therefore
\cite[Proposition~5.6]{BostChambertLoir2009}, applied with $D=P$, produces
$u_{p,\nu}$ with polar divisor a positive multiple of $P$ and
$U_{p,\nu}=\{|u_{p,\nu}|_p\le1\}$.  This is precisely the asserted rational
lemniscate description.  Notice that Proposition~5.6 is being applied
separately to each fixed domain; no continuity assertion is attributed to it.

On the subdivided reduction graph, the vertical correction of
\cite[Proposition~5.1]{BostChambertLoir2009} is the voltage generated by one
unit of current inserted at the vertex through $P$ and grounded at the cut
vertices.  Thus the coefficient at the source vertex is the effective
resistance of the parallel paths of lengths $\vartheta_\nu i(E)$.  By the
Eichler--Deuring mass formula \cite[\S5,\S10]{Deuring1941},
\[
 \sum_{[E]\ \mathrm{ss}}\frac1{i(E)}
 =2\sum_{[E]\ \mathrm{ss}}
   \frac1{\#\operatorname{Aut}(E)}
 =\frac{p-1}{12}.
\]
This proves the resistance formula in \eqref{eq:inner-resistance}.  If the
ramification index used for the subdivision is $e$, a retained branch of
normalized length $\vartheta_\nu i(E)$ consists of
$e\vartheta_\nu i(E)$ unit model edges, while each such edge contributes
$\log p/e$ to the normalized local degree.  Adjunction, as in
\eqref{eq:strict-cap-factor}, therefore gives
$\lambda_{p,\nu}=R_{p,\nu}\log p$, independently of the chosen
ramified extension, and proves the asserted limit.

For $p=2,3$, the explicit calculations in
\cite[Lemma~3.1 and the proof of Theorem~3.5]{Calegari2005} give
\[
 W_0(p)=\{|f_p|_p<p^{12/(p-1)}\}.
\]
The discs in \eqref{eq:wild-inner-discs} increase compactly to this domain.
Their complements are affinoid, and
\cite[Corollary~5.8]{BostChambertLoir2009}, together with
$f_p=q+O(q^2)$ and adjunction, gives the tangent gain displayed there.

For every $p$, the next member of the sequence supplies the larger compact
continuation neighborhood required in the finite-order estimates.  Finally,
the extensions above may depend on $\nu$.  Each prescribed finite local
extension can be induced by a finite number-field base change; then
\cref{lem:exact-base-change} shows that, after division by the extension
degree, local degrees, determinant degrees, intersection numbers, and slopes
inequalities are unchanged.  This proves the lemma.
\end{proof}

\begin{theorem}[Exact modular capacity pair at every prime]
\label{thm:allprime-capacity-pair}
For every prime $p$ and $0<Y<1/p$, the exhaustion of
\cref{lem:maximal-capacity-exhaustion}, with normalized invariants after the
harmless finite base changes used there, gives rational wide-opens
\[
 \Omega_{p,1}\Subset\Omega_{p,2}\Subset\cdots\Subset W_0(p)
\]
whose associated exact mixed continuation data $\alpha_\nu$ satisfy
\begin{align}
 \lim_{\nu\to\infty}\LambdaAH_{\alpha_\nu}
 &=L_p-2\pi Y,\label{eq:capacity-pair-Lambda}\\
 \lim_{\nu\to\infty}\Theta_{\alpha_\nu}
 &=L_p-2\pi Y,\label{eq:capacity-pair-Theta}\\
 \lim_{\nu\to\infty}\Energy_{\alpha_\nu}
 &=L_p-2\pi Y+C_p(Y),\label{eq:capacity-pair-Energy}
\end{align}
where
\[
 L_p=\frac{12\log p}{p-1}
\]
and $C_p(Y)$ is defined by \eqref{eq:Cp}.  Equivalently, the maximal modular
capacitary pair is
\begin{equation}\label{eq:capacity-pair-box}
 \boxed{
 (\Lambda_p(Y),\Energy_p(Y))
 =\left(
 \frac{12\log p}{p-1}-2\pi Y,
 \frac{12\log p}{p-1}-2\pi Y+C_p(Y)
 \right).}
\end{equation}
For every jointly rational packet $\mathcal S$ and every $1<\xi<d_pm$,
\begin{equation}\label{eq:allprime-explicit-slopes}
 \boxed{
 \tau^{\mathrm{curve}}_{p,\mathcal S}(\xi)
 \ge
 \left(1-\frac1m\right)
 \left(\frac{12\log p}{p-1}-2\pi Y\right)
 -\frac{C_p(Y)}m.}
\end{equation}
\end{theorem}

\begin{proof}
By \cref{lem:maximal-capacity-exhaustion}, the distinguished finite-place
contribution of the rational inner domains tends to
\[
 L_p=\frac{12\log p}{p-1}.
\]
At every finite place $\ell\ne p$ the chosen domain is the integral Tate
residue disc, so \cite[Example~5.11]{BostChambertLoir2009} contributes zero
relative tangent degree.  At infinity, \eqref{eq:allprime-qdisc} contributes
$-2\pi Y$.  This proves \eqref{eq:capacity-pair-Lambda}.

\smallskip
\noindent\emph{Target normal and Archimedean overflow.}
No interior point of the upper half-plane represents the cusp $P$.  Hence
\[
 \iota_Y^*P=o
\]
on the pointed disc.  The exact normal decomposition
\eqref{eq:theta-minus-lambda} therefore has no additional target-normal
collision, proving \eqref{eq:capacity-pair-Theta}.

It remains to evaluate the Archimedean term $\mathscr W$.  On the boundary
write $\tau=x+iY$, $0\le x<1$.  Nonidentity self-collisions of the modular
$q$-disc are parametrized by primitive bottom rows $(c,d)$ with
$c>0$ and $p\mid c$.  For
$\gamma=\left(\begin{smallmatrix}a&b\\c&d\end{smallmatrix}\right)$,
\[
 \Im(\gamma\tau)=\frac{Y}{(cx+d)^2+c^2Y^2}.
\]
The source Green function is $2\pi(\Im\tau-Y)$, so the direct-image
Green--Jensen formula gives
\begin{equation}\label{eq:allprime-collision-unfold}
 \mathscr W_p(Y)
 =\sum_{\substack{c>0\\p\mid c}}
   \sum_{\substack{d\in\Z\\(c,d)=1}}
   \int_0^1 2\pi Y
   \left(\frac1{(cx+d)^2+c^2Y^2}-1\right)_+\dd x.
\end{equation}
Elliptic stabilizers occur only over a measure-zero set and do not alter this
integral.  For fixed $c$, the summand vanishes when $cY\ge1$.  When $cY<1$,
unfolding the primitive residues $d\bmod c$ and putting
$b_c=\sqrt{1-c^2Y^2}$ gives
\begin{align*}
 &\frac{2\pi Y\varphi(c)}c
 \int_{-b_c}^{b_c}
 \left(\frac1{u^2+c^2Y^2}-1\right)\dd u\\
 &\qquad=
 \frac{4\pi\varphi(c)}{c^2}
 \left(\arccos(cY)-cY\sqrt{1-c^2Y^2}\right).
\end{align*}
Thus
\begin{equation}\label{eq:allprime-W-C}
 \mathscr W_p(Y)=C_p(Y).
\end{equation}
Together with \eqref{eq:energy-theta}, this proves
\eqref{eq:capacity-pair-Energy}.

The rational skeleton exhaustion of
\cref{lem:maximal-capacity-exhaustion} supplies the required inner data.
Use the same Archimedean disc for every member.  Its distinguished finite
contribution tends to $L_p$, while the target-normal comparison has no extra
cusp preimage and the Archimedean term is the fixed number $C_p(Y)$.
Consequently the exact mixed invariants converge to
\eqref{eq:capacity-pair-Lambda}--\eqref{eq:capacity-pair-Energy}.  Since
$\log p>1-1/p$, one has
$L_p>12/p>2\pi/p>2\pi Y$; hence every sufficiently large member is
pseudoconcave.  Apply \cref{thm:curve-Bost-exact} to those members and pass
to the limit.  The arithmetic type is independent of the continuation datum,
so this yields \eqref{eq:allprime-explicit-slopes}.
\end{proof}

\begin{remark}[Why genus zero was inessential for the collision formula]
The proof of \eqref{eq:allprime-W-C} uses only the uniformizing quotient
$\Gamma_0(p)\backslash\mathfrak H$, the width-one cusp, and the source Green
function.  A Hauptmodul is unnecessary.  Consequently the formula
$\Energy_p(Y)=\Lambda_p(Y)+C_p(Y)$ from \cref{prop:Jensen} is intrinsic and
extends to every prime level.
\end{remark}

\section{All-prime density one}
\label{sec:allprime-density}

\begin{proof}[Proof of \cref{thm:density-all-primes}]
Fix $\beta>1$.  Let $\mathcal S_j$ be the rational odd arguments in the shell
\[
 A=\beta^j\le s<\beta A,
\]
and put $k_j=\#\mathcal S_j$.  Empty shells may be ignored.  For a nonempty
shell,
\[
 k_jA\le m-1<k_j\beta A,
 \qquad s_*<\beta A,
 \qquad Q_{\mathcal S_j}<\beta A m.
\]
Choose $\xi=1+1/\log A$.  Since $d_p$ is fixed,
\[
 G_{d_p,m}(\xi)=d_p\log m+O_p(1),
\]
and \eqref{eq:curve-tau} gives
\begin{equation}\label{eq:allprime-shell-upper}
 \tau^{\mathrm{curve}}_{p,\mathcal S_j}(\xi)
 \le\frac{2d_p\beta}{k_j}\log m
 +O_{p,\beta}\!\left(\frac1{k_j}\right).
\end{equation}

In \cref{thm:allprime-capacity-pair} choose $Y=1/m$; for a nonempty shell
this is $<1/p$ once $j$ is sufficiently large.  Since
\[
 0\le\arccos x-x\sqrt{1-x^2}\le\frac\pi2,
 \qquad \varphi(c)\le c,
\]
we have
\begin{equation}\label{eq:allprime-C-bound}
 C_p(1/m)
 \le\frac{2\pi^2}{p}
 H_{\floor{(m-1)/p}}
 =O_p(\log m).
\end{equation}
The explicit exact inequality \eqref{eq:allprime-explicit-slopes} therefore
gives
\begin{equation}\label{eq:allprime-shell-lower}
 \tau^{\mathrm{curve}}_{p,\mathcal S_j}(\xi)
 \ge L_p-O_p\!\left(\frac{\log m}{m}\right),
 \qquad
 L_p=\frac{12\log p}{p-1}.
\end{equation}

Initially $k_j=O_\beta(A)$ and hence $m=O_\beta(A^2)$, so
$\log m=O_\beta(\log A)$.  The right side of
\eqref{eq:allprime-shell-lower} is at least $L_p/2$ for all sufficiently large
$j$.  Combining the upper and lower estimates first gives
\[
 k_j=O_{p,\beta}(\log A).
\]
It follows that
\[
 m=O_{p,\beta}(A\log A),
 \qquad
 \log m=\log A+O_{p,\beta}(\log\log A).
\]
Substitution back yields
\begin{equation}\label{eq:allprime-shell-count}
 k_j\le
 \left(\frac{2d_p\beta}{L_p}+o_p(1)\right)\log A.
\end{equation}

If $\beta^J\le X<\beta^{J+1}$, summing over $j\le J$ gives
\begin{equation}\label{eq:allprime-explicit-constant}
 R_p(X)\le
 \left(\frac{d_p\beta}{L_p\log\beta}+o_p(1)\right)(\log X)^2.
\end{equation}
The function $\beta/\log\beta$ is minimized at $\beta=e$, so
\[
 \frac{ed_p}{L_p}=\frac{e\,d_p(p-1)}{12\log p},
\]
which proves \eqref{eq:density-all-primes}.  When $g_p\ge1$, substitute
$d_p\le2g_p$ from \cref{thm:finite-map-zero} to obtain
\eqref{eq:density-all-primes-genus}.

For each shell, the packet is fixed before the source-degree limit
$D\to\infty$ is taken.  Packet-dependent module generators, bad primes,
metric comparisons, and boundary norms have therefore disappeared into
$o(D^2)$ before $j\to\infty$.  The subsequent limits in the capacitary
exhaustion and in $j$ introduce no packet-uniform auxiliary-prime assumption.
\end{proof}

\begin{remark}[What the all-prime theorem does not assert]
The coefficient in \eqref{eq:density-all-primes} is explicit for any chosen
finite map of degree $d_p$, but the general construction
$d_p\le2g_p$ need not be gonality-optimal.  Improving the map can therefore
improve the displayed constant and the finite-packet bounds below.  No
uniformity as $p\to\infty$ beyond the stated formula is claimed, and the
theorem still permits an infinite polylogarithmically sparse set of rational
exceptions.
\end{remark}



\subsection{Certified finite packets at positive genus}
\label{sec:positive-genus-finite}

The exact capacity formula permits a finite-packet certificate that does not
enumerate subsets.

\begin{proposition}[One-dimensional finite-packet certificate]
\label{prop:finite-packet-certificate}
Fix a prime $p$, let $d_p$ be an admissible finite-map degree in
\cref{thm:finite-map-zero}, and let $\bar d\ge d_p$ be a certified upper
bound.  Fix an odd cutoff $S$ and an integer $k\ge1$.  Put
\begin{align}
 m_0(k,x)&=x+k^2,\label{eq:finite-cert-m0}\\
 U_{\bar d,k}(x)&=2\bar d\left[
 \frac{x-2+x\log m_0(k,x)}{m_0(k,x)}
 +\frac{x+2}{m_0(k,x)^2}\right],\label{eq:finite-cert-U}\\
 B_{p,k}(x)&=
 \left(1-\frac1{m_0(k,x)}\right)
 \left(L_p-\frac{2\pi}{m_0(k,x)}\right)
 -\frac{2\pi^2}{p\,m_0(k,x)}
 H_{\floor{(m_0(k,x)-1)/p}}.\label{eq:finite-cert-B}
\end{align}
Assume, for every relevant $x$, that
\[
 m_0(k,x)>\max\{p,\bar d\},
 \qquad L_p-\frac{2\pi}{m_0(k,x)}>0.
\]
If
\begin{equation}\label{eq:finite-cert-condition}
 B_{p,k}(x)>U_{\bar d,k}(x)
 \qquad(2k+1\le x\le S,\ x\text{ odd}),
\end{equation}
then among the odd arguments $3,5,\ldots,S$ at most $k-1$ values
$\zeta_p(s)$ are rational.
\end{proposition}

\begin{proof}
Suppose instead that $k$ distinct values are rational and let $\mathcal S$ be
one such $k$-element packet, with $x=\max\mathcal S$.  The other $k-1$
arguments are distinct positive odd integers at least $3$, so their sum is at
least
\[
 3+5+\cdots+(2k-1)=k^2-1.
\]
Thus
\begin{equation}\label{eq:finite-cert-m-lower}
 m=1+\sum_{s\in\mathcal S}s\ge x+k^2=m_0(k,x).
\end{equation}
For every $s\in\mathcal S\setminus\{x\}$ one has $s\le x-2$, and hence
$s^2\le(x-2)s$.  Therefore
\begin{equation}\label{eq:finite-cert-Q}
 Q_{\mathcal S}
 \le x^2+(x-2)(m-1-x)
 =(x-2)m+x+2.
\end{equation}
Choose $\xi=\bar d$.  The assumption $\bar d<m_0(k,x)\le m$ makes
this cutoff admissible.  Since $\bar d\ge d_p$, the second branch of
\eqref{eq:Gdm} gives
\[
 G_{d_p,m}(\bar d)
 =d_p\log\!\left(\frac{d_pm}{\bar d}\right)
 \le d_p\log m\le\bar d\log m.
\]
Equations \eqref{eq:curve-tau} and \eqref{eq:finite-cert-Q} therefore give
\[
 \tau^{\mathrm{curve}}_{p,\mathcal S}(\bar d)
 \le2\bar d\left[
 \frac{x-2+x\log m}{m}+\frac{x+2}{m^2}\right].
\]
The derivative of the quantity in square brackets is
\[
 \frac{2-x\log m}{m^2}-\frac{2(x+2)}{m^3}<0
\]
for $x\ge3$ and $m\ge3$.  Hence \eqref{eq:finite-cert-m-lower} implies
\begin{equation}\label{eq:finite-cert-tau-upper}
 \tau^{\mathrm{curve}}_{p,\mathcal S}(\bar d)\le U_{\bar d,k}(x).
\end{equation}

On the analytic side use \cref{thm:allprime-capacity-pair} with the fixed
choice $Y=1/m_0(k,x)$.  Write
\[
 A=L_p-\frac{2\pi}{m_0(k,x)}.
\]
For fixed $Y$, the exact lower side is
\[
 \left(1-\frac1m\right)A-\frac{C_p(Y)}m
 =A-\frac{A+C_p(Y)}m,
\]
which is increasing in $m$ because $A>0$ and $C_p(Y)\ge0$.  Moreover
\eqref{eq:allprime-C-bound}, with $m_0$ in place of $m$, gives
\[
 C_p(1/m_0)
 \le\frac{2\pi^2}{p}
 H_{\floor{(m_0-1)/p}}.
\]
Consequently joint rationality forces
\begin{equation}\label{eq:finite-cert-analytic-lower}
 \tau^{\mathrm{curve}}_{p,\mathcal S}(\bar d)\ge B_{p,k}(x).
\end{equation}
Equations \eqref{eq:finite-cert-tau-upper} and
\eqref{eq:finite-cert-analytic-lower} contradict
\eqref{eq:finite-cert-condition}.
\end{proof}

\begin{proof}[Proof of \cref{thm:positive-genus-counts}]
Apply \cref{prop:finite-packet-certificate} with $S=2401$, $\bar d=2g_p$, and the values
shown below.  The certificate checks every admissible odd largest argument
$x$, using outward-rounded integer fixed-point intervals at scale $10^{45}$.
The harmonic numbers in \eqref{eq:finite-cert-B} are enclosed by summing the
individual rational terms outward; $\pi$ is enclosed by Machin's formula and
alternating arctangent series, and logarithms are enclosed after dyadic range
reduction by the positive $\operatorname{arctanh}$ series.  No binary
floating-point value enters a certificate decision.

\begin{center}
\begin{tabular}{@{}crrrrr@{}}
\toprule
$p$&$\bar d$&excluded $k$&rational cap&worst $x$&certified margin\\
\midrule
$11$&$2$&$191$&$190$&$2401$&$>0.019624155425482$\\
$17$&$2$&$227$&$226$&$2401$&$>0.006340667081541$\\
$19$&$2$&$238$&$237$&$2401$&$>0.013201735527294$\\
$23$&$4$&$378$&$377$&$2401$&$>0.006562285899278$\\
$29$&$4$&$415$&$414$&$2401$&$>0.005457529203788$\\
$31$&$4$&$426$&$425$&$2401$&$>0.002866264645707$\\
$37$&$4$&$458$&$457$&$2401$&$>0.002720591459553$\\
$41$&$6$&$596$&$595$&$2401$&$>0.003315013860040$\\
$43$&$6$&$607$&$606$&$2401$&$>0.000684777792175$\\
\bottomrule
\end{tabular}
\end{center}
Subtracting the rational caps from $1200$ gives the irrational counts in
\cref{thm:positive-genus-counts}.

For the final column of that theorem, if a prefix contains $N$ odd arguments,
we apply the same certificate with $k=N-999$: excluding such a packet leaves
at most $N-1000$ rational values.  The certified cutoffs are respectively
\[
 2379,\ 2459,\ 2483,\ 2823,\ 2921,\ 2951,\ 3039,\ 3445,\ 3483.
\]
The stored certificate also checks the immediately preceding cutoff for this
sufficient criterion and records a nonpositive lower margin there.  The
certificate files are listed in \cref{sec:certificate-files}.
\end{proof}

\section{Logical dependency and adversarial audit}\label{sec:audit}

\begingroup\small
\begin{longtable}{@{}>{\raggedright\arraybackslash}p{0.20\textwidth}>{\raggedright\arraybackslash}p{0.33\textwidth}>{\raggedright\arraybackslash}p{0.39\textwidth}@{}}
\toprule
Gate&Input&Use and final audit conclusion\\
\midrule
Single blocks&Calegari; BGG/Bol; internal \cref{sec:blocks}&Constructs each algebraic extension and exact raw rows for general $p$.\\
Joint rank and descent&Internal \cref{thm:joint-rank,lem:allprime-descent}&Even symmetric powers descend away from the fixed elliptic divisor; the multiplicity-free direct sum has full monodromy rank.\\
Genus-zero zero estimate&CDT Theorem~3.2.13; \cref{cor:joint-zero}&Applies on $\mathbf P^1$ and retains the sharp normalized interval $m+o(1)$.\\
General-curve zero estimate&CDT Theorem~3.2.8; \cref{thm:finite-map-zero}&Restriction of scalars gives a derivative-stable rank-$d_pm$ system; $d_p$ enlarges only the pivot range, not the genuine source rank.\\
Hasse saving&Katz; \cref{prop:single-action,thm:block-Hasse}&Used only for the sharp genus-zero theorems; not load-bearing for arbitrary-$p$ density one.\\
Genus-zero Cartier&Exact Dwork; \cref{lem:no-backflow,prop:multi-Cartier}&Common Hauptmodul Frobenius and full-product digitization give the CDT degree-$D$ window.\\
General-curve Cartier&Exact $q$-Dwork; \cref{lem:q-jet-injective,lem:q-block,prop:q-jet-Cartier}&Global divisor control makes the short $q$-jet injective; exact $q\mapsto q^\ell$ removes all Taylor correction terms.\\
Fitting polygon&Internal \cref{prop:Fitting,prop:capped-window}&Counts $\ell$-adic thickness $s-b+1$; used for strengthening, not for the certified finite counts.\\
Genus-zero Bost bridge&Bost, Bost--Charles, CDT; \cref{lem:fixed-bundle}&The $\mathbf P^1$ fixed-bundle comparison changes only $o(D^2)$.\\
General-curve fixed-twist bound&Finite generation; Bost--Chambert-Loir Theorem~6.1, (6.7)--(6.8); \cref{lem:robust-fixed-twist}&A fixed packet is a quotient of finitely many fixed scalar twists.  Relative Serre vanishing identifies the quotient lattices with the good-prime coefficient lattices.  Packet dependence is $O(\log D)$ per pivot, while the coefficient of $D$ depends only on $p$.\\
Scalar capacitary jets&Bost--Chambert-Loir Proposition~5.15; Cauchy; \cref{lem:robust-cap-jet}&Contraction is performed first, leaving a scalar section.  The arbitrary reference tangent norm cancels formally and gives the term $-n\Lambda_p^\circ$.\\
Hybrid compatibility&Raw depth; intrinsic Cartier; \cref{prop:curve-hybrid-compatibility}&Shrinking the finite source lattice improves analytic operator norms.  Small-prime residuals require proper prime powers; coherent models eliminate any hidden packet-dependent $D^2$ term at fixed bad primes.  Large primes give the stated window integral.\\
$p=13$ continuation&Maier; Deuring; Buzzard; \cref{lem:p13-annulus}&Identifies $W_0(13)=\{|f_{13}|_{13}<13\}$ and permits explicit finite estimates.\\
All-prime positive width&Buzzard; Bost--Chambert-Loir Proposition~5.1, the proof of Proposition~5.6, Example~5.11, and Proposition~5.15; \cref{lem:strict-cap-collar,prop:positive-width}&The capacitary vertical coefficient on the component met by the cusp is strictly positive, so adjunction makes the local tangent norm strictly smaller than the integral-disc norm; a complex $q$-disc consumes arbitrarily little of that gain.\\
Exact mixed capacitary normalization&Bost--Chambert-Loir Propositions~5.1, 5.6, and~5.17, and Corollary~5.20; Bost--Charles Lemma~7.2.7, Proposition~7.2.10, and the meromorphicity criterion of \S7.3; Chen--Moriwaki Corollary~4.5.2; \cref{thm:curve-Bost-exact,app:exact-normalization}&The finite capacitary $\Q$-model and Archimedean direct-image Green functions define one nef and big mixed divisor.  Denominator clearing realizes the capacitary tangent line as a pseudoconcave normal line, but neither a whole-divisor direct image nor a meromorphic map is inferred from that normal-line isometry.  The exact source formula uses the Chen--Moriwaki determinant theorem; Bost--Charles Theorem~8.2.5 is a theta formula and is not used as a determinant formula.  A nested lemniscate gives the uniform finite-order estimate and the exact target-normal $\Theta$-cancellation.\\
Exact modular capacity pair&Buzzard Proposition~4.1; Coleman--McMurdy; Eichler--Deuring mass formula; Bost--Chambert-Loir Propositions~5.1 and~5.6; \cref{lem:maximal-capacity-exhaustion,thm:allprime-capacity-pair}&After an unramified splitting extension, the supersingular annuli for $p\ge5$ are parallel resistors of widths $\#\operatorname{Aut}(E)/2$.  Rational skeleton cuts give inner wide-opens whose effective resistances converge to $12/(p-1)$; Proposition~5.6 is used only to realize each cut-domain as a rational lemniscate.  The wild primes $2,3$ use their explicit Hauptmodul discs.  The Archimedean collision sum is intrinsic to the $\Gamma_0(p)$ quotient and is genus-independent.\\
Quantitative genus-zero density&Internal \cref{sec:density}&Uses exact modular width and collision formula, yielding the coefficient $e(p-1)/(12\log p)$.\\
All-prime density&Internal \cref{thm:allprime-capacity-pair,sec:allprime-density}&Uses the exact capacity pair with $Y=1/m$ and multiplicative shells, yielding the explicit coefficient $e d_p(p-1)/(12\log p)$.  The robust theorem remains an independent qualitative proof.\\
Numerics&Independent fixed-point certificates; \cref{prop:p13-finite,prop:finite-packet-certificate}&The genus-zero counts retain their original certificates.  The new positive-genus table uses the certified upper bound $\bar d=2g_p$, checks every possible largest packet argument with outward-rounded integer fixed-point arithmetic, and uses no binary floating point.\\
\bottomrule
\end{longtable}
\endgroup

\begin{warning}[Distinct weights]
The functional-rank proof uses that the representations
$\Sym^{s-1}(\C^2)$ occur with multiplicity one.  Repeated copies of the same
weight can support diagonal subrepresentations and do not automatically add
functional rank.
\end{warning}

\begin{warning}[No scalarization]
The proof does not construct global scalar coordinates carrying literal depths
$0,1,\ldots,s$.  The small-prime saving uses only the mod-$\ell$ Hasse
comparison of the genuine deepest multiplier.  The large-prime argument stays
on the de Rham frame torsor and digitizes the complete truncated products.
\end{warning}

\begin{warning}[Counting versus individual identification]
\Cref{thm:counts,thm:positive-genus-counts} bound the number of rational
values in finite ranges but
does not identify every possible exception.  Proving long explicit lists of
individual irrational values requires a different amplification: an
unconditional arithmetic anchor packet that contributes functional rank under
the rationality assumption for one target only.
\end{warning}

\begin{warning}[All-prime scope]
The all-prime theorem is a fixed-$p$ assertion.  Its leading density
coefficient is now explicit once a finite map degree $d_p$ is chosen, and the
safe construction $d_p\le2g_p$ is uniform at fixed level.  The coefficient is
not claimed optimal: the displayed choice $2g_p$ is only an upper bound for the actual degree, and a lower-degree map unramified at the Tate cusp improves
both the density constant and the finite-packet arithmetic cost.  The sharp
Hasse-improved type and denominator polygon remain genus-zero results, whereas
the exact modular capacity and collision formulas are valid at every prime.
\end{warning}

\begin{warning}[Exact-normalization scope]
The normalization-robust theorem still supplies an independent qualitative
proof of density one, but the explicit leading coefficient in
\cref{thm:density-all-primes} uses the exact mixed normalization and the
capacity exhaustion of \cref{thm:allprime-capacity-pair}.  The exact result
\cref{thm:curve-Bost-exact} realizes the metrized tangent line, not the entire
finite vertical Bost--Chambert-Loir divisor, as the normal line of a
formal-analytic surface.  It also does not infer a Bost--Charles meromorphic
map from an arbitrary identification of generic formal discs.  The finite
divisor is used directly in the mixed canonical metric, and the determinant
source formula is supplied by Chen--Moriwaki rather than by the Bost--Charles
theta Hilbert--Samuel theorem.  These distinctions are recorded explicitly in
\cref{warn:normal-line-only}.
\end{warning}

\appendix


\section{Exact mixed capacitary normalization and the normal-line bridge}
\label{app:exact-normalization}

This appendix proves \cref{thm:curve-Bost-exact}.  The construction is
adelic but deliberately asymmetric.  At finite places we use the
Bost--Chambert-Loir capacitary $\Q$-model itself.  At Archimedean places we
use the direct-image Green function of the chosen continuation map in the
sense of Bost--Charles.  The first audit point is that these two pieces define
one semipositive adelic divisor on the modular curve.  The second is that the
capacitary tangent line, after denominator clearing, is a genuine normal line
of a pseudoconcave formal-analytic surface.

No compatible Bost--Charles meromorphic map from that formal-analytic surface
to the modular curve is asserted or needed.  Such a map is a separate
fraction-field condition in the formal category and does not follow merely
from an identification of generic formal discs.  Likewise, an isometry of
normal lines does not identify the full finite vertical divisor.  The exact
source and jet formulas below are therefore proved directly on $X_p$, rather
than by invoking an unproved whole-divisor direct-image realization.

\subsection{The canonical mixed capacitary divisor}

Fix $p$ and the continuation datum
\[
 \boldsymbol\alpha=
 ((\Omega_v,o_v,\iota_v))_v=\boldsymbol\Omega_p^\circ
\]
of \cref{prop:positive-width}.  At a finite place, $\iota_v$ is the inclusion
of a connected wide-open neighborhood of $P$ in $X_p^{\mathrm{an}}$; at all
but finitely many places it is the integral Tate residue disc.  At an
Archimedean place $\sigma$, write
\[
 g_{\sigma}=g_{\Omega_\sigma,o_\sigma},\qquad
 \mu_\sigma=\mu_{\Omega_\sigma,o_\sigma}
\]
for the equilibrium Green function and harmonic measure.  The map
$\iota_\sigma$ is nonconstant and unramified at $o_\sigma$.

At the finite places, \cite[Propositions~5.1 and~5.17]
{BostChambertLoir2009} supplies a normal flat projective model
$\mathscr X_\alpha$ of $X_p$ and a vertical $\Q$-divisor
$V_\alpha^{\mathrm{fin}}$ such that
\[
 \mathscr P_\alpha^{\mathrm{fin}}
 =\mathscr P+V_\alpha^{\mathrm{fin}}
\]
induces the capacitary metric of $\Omega_v$ on $\mathcal O(P)$ at every finite
place.  After one fixed birational refinement, which does not change these
metrics or their intersection numbers, we may and do assume that the section
$P$ lies in the regular locus of $\mathscr X_\alpha$; thus adjunction and
$P^*\mathcal O(P)$ are literal near the section.  The adelic tube is defined
using the integral Tate cusp model at every
finite place other than $p$ and an adapted model of the chosen rational
wide-open at $p$.  Proposition~5.17, through its Raynaud comparison and
Moret--Bailly descent step, places these local data on one global model.  The
corresponding vertical correction is zero at every integral Tate place.  In
\cref{prop:exact-hybrid-assembly} the coefficient-bundle models used in
\cref{sec:qjet} are chosen from this same model after one common refinement
and enlargement of the fixed exceptional set.  The finite vertical divisor is
independent of the auxiliary Archimedean components by the local uniqueness in
Proposition~5.1.

For an Archimedean embedding $\sigma$, let
\begin{equation}\label{eq:direct-image-green}
 g_{\alpha,\sigma}=\iota_{\sigma*}g_\sigma
\end{equation}
be the direct-image Green function for $P$ in the sense of
\cite[Sections~5.1 and~7.2]{BostCharles2022}.  It has logarithmic singularity
of multiplicity one at $P$, because $\iota_\sigma$ is unramified at
$o_\sigma$.  Its curvature measure is the positive direct image of the
source equilibrium measure.

\begin{definition}[Canonical mixed divisor]\label{def:mixed-divisor}
Define
\begin{equation}\label{eq:mixed-divisor}
 \widehat P_\alpha=
 \left(\mathscr P_\alpha^{\mathrm{fin}},
       (g_{\alpha,\sigma})_{\sigma}\right)
 \in\widehat{\operatorname{Div}}(\mathscr X_\alpha)_\Q
\end{equation}
and put
\[
 \overline L_\alpha=\mathcal O(\widehat P_\alpha).
\]
Let $\overline T_\alpha$ denote the line $T_PX_p$ with the finite
Bost--Chambert-Loir capacitary metrics and, at infinity, the source
capacitary metric on $T_{o_\sigma}\Omega_\sigma$ transported through
$D\iota_\sigma(o_\sigma)$.  Finally put
\begin{equation}\label{eq:mixed-invariants}
 \LambdaAH_\alpha=\widehatdeg\overline T_\alpha,
 \qquad
 \Theta_\alpha=\widehatdeg P^*\overline L_\alpha,
 \qquad
 \Energy_\alpha=\widehat P_\alpha^{\,2}.
\end{equation}
\end{definition}

The distinction between $\LambdaAH_\alpha$ and $\Theta_\alpha$ is essential.
The former is a source-tangent degree.  The latter is the degree of the target
normal line determined by the direct-image Green function.  Other preimages
of $P$ under $\iota_\sigma$ may make the two different.

\begin{lemma}[Base-change and denominator invariance]
\label{lem:exact-base-change}
Let $K/\Q$ be a finite extension.  Pull back all models, $\Q$-divisors,
metrized lines, sources, and jet targets to $\mathcal O_K$.  Then every
Arakelov degree and arithmetic intersection number is multiplied by
$[K:\Q]$.  Consequently all identities and inequalities in this appendix are
unchanged after division by $[K:\Q]$.  The same conclusions hold for an
Arakelov $\Q$-divisor after multiplying by a common denominator and dividing
the resulting degree or intersection number by the corresponding first or
second power.
\end{lemma}

\begin{proof}
The local degree at a place $v$ is repeated with the usual local-degree
weights over the places of $K$ above $v$; their sum is $[K:\Q]$.  This proves
the assertion for metrized lines and for determinant degrees.  The projection
formula gives the same scaling for arithmetic intersections.  Homogeneity in
a common denominator gives the final assertion.
\end{proof}

\begin{proposition}[Nefness, bigness, and the exact normal decomposition]
\label{prop:exact-mixed-divisor}
The divisor $\widehat P_\alpha$ is effective, nef, and big.  Moreover
\begin{equation}\label{eq:theta-minus-lambda}
 \Theta_\alpha-\LambdaAH_\alpha
 =\sum_{\sigma}
   \int_{\Omega_\sigma}g_\sigma\,
   \delta_{\iota_\sigma^*P-o_\sigma}\ge0
\end{equation}
and
\begin{equation}\label{eq:energy-theta}
 \boxed{
 \Energy_\alpha=\Theta_\alpha+\mathscr W_\alpha,\qquad
 \mathscr W_\alpha=
 \sum_\sigma\int_{X_{p,\sigma}(\C)}
 g_{\alpha,\sigma}\,c_1(\overline L_{\alpha,\sigma})
 =\sum_\sigma\int_{\Omega_\sigma}
 g_\sigma\,\iota_\sigma^*c_1(\overline L_{\alpha,\sigma})\ge0.}
\end{equation}
In particular,
\[
 \Theta_\alpha\ge\LambdaAH_\alpha>0,
 \qquad
 \Energy_\alpha>0.
\]
The finite capacitary corrections occur in $\Theta_\alpha$ exactly once; the
term $\mathscr W_\alpha$ is purely Archimedean.
\end{proposition}

\begin{proof}
At each finite place, the vertical correction in
$\mathscr P_\alpha^{\mathrm{fin}}$ is the solution of the graph Dirichlet
problem of \cite[Proposition~5.1]{BostChambertLoir2009}.  The resulting
$\Q$-divisor is effective.  It has intersection zero with every vertical
component occurring in its support and nonnegative intersection with every
other vertical component.  At infinity, the direct-image Green function is
nonnegative and its curvature is the positive direct image of the source
equilibrium measure.  Hence $\widehat P_\alpha$ is effective and defines a
semipositive adelic metric on $\mathcal O(P)$.

At a finite place the continuation map is the inclusion, so the source and
target capacitary tangent metrics agree.  At an Archimedean place, the local
normal comparison in the proof of
\cite[Lemma~7.2.7 and Proposition~7.2.10]{BostCharles2022}, with ramification
index one, gives
\[
 \widehatdeg_\sigma P^*\overline L_\alpha
 -\widehatdeg_\sigma\overline T_\alpha
 =\int_{\Omega_\sigma}g_\sigma\,
   \delta_{\iota_\sigma^*P-o_\sigma}.
\]
The divisor $\iota_\sigma^*P-o_\sigma$ is effective.  Summing proves
\eqref{eq:theta-minus-lambda}, so the horizontal component $P$ has strictly
positive arithmetic degree.  The numerical criterion of
\cite[Proposition~6.9]{Bost1999}, used in
\cite[Proposition~7.7]{BostChambertLoir2009}, now proves nefness.  Indeed, the
curvature is nonnegative, the degree on $P$ is positive, the degree on every
vertical component in the support is zero, and the remaining vertical degrees
are nonnegative; every horizontal curve different from $P$ has nonnegative
degree because both the finite divisor and the Green functions are
effective.

Write $\mathscr P_\alpha^{\mathrm{fin}}=\mathscr P+V$.  By the orthogonality
in \cite[Proposition~5.1]{BostChambertLoir2009},
$\mathscr P_\alpha^{\mathrm{fin}}\cdot V=0$.  The induction formula for
the arithmetic intersection pairing, in the normalization used in the proof
of \cite[Proposition~7.2.6]{BostCharles2022}, therefore gives
\[
 \widehat P_\alpha^{\,2}
 =\widehatdeg P^*\overline L_\alpha
  +\sum_\sigma\int_{X_{p,\sigma}(\C)}
    g_{\alpha,\sigma}\,c_1(\overline L_{\alpha,\sigma}).
\]
The projection formula for direct-image Green functions gives the second
expression in \eqref{eq:energy-theta}.  Both factors in its integrand are
nonnegative.  Thus the self-intersection is positive.  We have already proved
that $\overline L_\alpha$ is nef; a nef Arakelov divisor of positive
self-intersection is nef and big in the sense of
\cite[Section~8.2.2]{BostCharles2022}.  This proves the proposition.
\end{proof}

\subsection{The capacitary normal-line realization}

The preceding proposition already constructs the exact target metric needed
for arithmetic Hilbert--Samuel.  The following result closes the separate
formal-analytic realization issue.  Its statement is intentionally about the
normal line, not the entire finite vertical divisor.

\begin{proposition}[Denominator-cleared capacitary normal realization]
\label{prop:cap-normal-realization}
After a finite extension $K/\Q$ clearing the rational finite-place exponents
of $\overline T_\alpha$, there exists a smooth formal-analytic arithmetic
surface
\[
 \widetilde{\mathscr V}_\alpha
 =\left(\widehat{\mathscr V}_\alpha,
       (\Omega_\sigma,o_\sigma,\gamma_\sigma)_\sigma\right)
 \quad\text{over }\mathcal O_K
\]
whose metrized normal line along its canonical section is isometric to
$\overline T_\alpha\otimes_\Q K$.  Consequently
\begin{equation}\label{eq:normal-realization-degree}
 \frac1{[K:\Q]}
 \widehatdeg N_{P_K}\widetilde{\mathscr V}_\alpha
 =\LambdaAH_\alpha>0,
\end{equation}
so $\widetilde{\mathscr V}_\alpha$ is pseudoconcave.
\end{proposition}

\begin{proof}
By \cite[Propositions~5.1 and~5.17]{BostChambertLoir2009}, the finite metrics
on $T_PX_p$ are $\Q$-model metrics and are integral at all but finitely many
places.  Choose a common denominator for their valuation exponents and a
finite extension $K/\Q$ whose ramification indices above the exceptional
places are divisible by that denominator; one may take a compositum of
suitable Eisenstein extensions at the finitely many rational primes involved.
The unit balls of the pulled-back
metrics are then honest rank-one $\mathcal O_{K,w}$-lattices.  Since only
finitely many differ from the standard tangent lattice, they are the
localizations of one invertible fractional $\mathcal O_K$-module
$\mathcal N_\alpha\subset T_PX_p\otimes_\Q K$.

Set
\[
 \widehat{\mathscr V}_\alpha
 =\operatorname{Spf}
   \widehat{\operatorname{Sym}}_{\mathcal O_K}
   (\mathcal N_\alpha^\vee).
\]
This is a smooth one-dimensional formal scheme over $\mathcal O_K$, with
zero section $P_K\simeq\operatorname{Spec}\mathcal O_K$ and normal module
$\mathcal N_\alpha$.  For each embedding $\sigma:K\hookrightarrow\C$, choose
a formal coordinate at $o_\sigma$ and a conjugation-compatible formal
isomorphism
\[
 \gamma_\sigma:
 \widehat{\mathscr V}_{\alpha,\sigma}
 \xrightarrow{\sim}\widehat{\Omega}_{\sigma,o_\sigma}
\]
whose differential identifies the complex fibre of $\mathcal N_\alpha$ with
$T_{o_\sigma}\Omega_\sigma$ after transport through
$D\iota_\sigma(o_\sigma)$.  In dimension one the derivative may be prescribed
by multiplying a local coordinate by a nonzero complex scalar; the choices
may be paired under complex conjugation.  Equip the complex fibre of
$\mathcal N_\alpha$ with the corresponding source capacitary norm.  The data
then form a smooth formal-analytic arithmetic surface in the sense of
\cite[Sections~6.1--6.2]{BostCharles2022}, and its metrized normal line is
precisely $\overline T_\alpha\otimes_\Q K$.

Taking normalized Arakelov degrees and applying
\cref{lem:exact-base-change} gives
\eqref{eq:normal-realization-degree}.  Positivity is
\cref{prop:positive-width}; by definition, it is the pseudoconcavity of
$\widetilde{\mathscr V}_\alpha$.
\end{proof}

\begin{warning}[Normal-line realization is not a map]
\label{warn:normal-line-only}
An isometry of metrized normal lines determines neither a vertical divisor on
the whole special fibre nor a Bost--Charles meromorphic map to
$\mathscr X_{\alpha,K}$.  In the formal category, meromorphicity requires the
coordinate functions of the generic map to lie in the fraction field of the
formal function ring and to become regular after a modification; an arbitrary
identification of generic formal discs does not establish this condition.
Accordingly, \cref{prop:cap-normal-realization} makes no such assertion.  The
exact normalization below uses the Bost--Chambert-Loir finite divisor itself
and is proved directly on the modular curve.
\end{warning}

\subsection{Exact arithmetic Hilbert--Samuel for the source}

Regard $\overline L_\alpha$ as a semipositive adelic line bundle on the
smooth projective curve $X_p/\Q$.  Its underlying line bundle
$L=\mathcal O(P)$ has degree one and is therefore ample.  Give the fixed
bundle $\calF_{\mathcal S}$ any fixed dominated adelic metric compatible with
its integral models.  On spaces of sections use the induced supremum norms;
at every Archimedean place replace the supremum unit ball by its John
ellipsoid in order to obtain a Hermitian norm.

\begin{proposition}[Exact determinant source formula]
\label{prop:exact-source-formula}
For every fixed packet $\mathcal S$,
\begin{equation}\label{eq:appendix-exact-source}
 \widehatdeg\overline M_D^\alpha
 =\frac m2\Energy_\alpha D^2+o_{p,\mathcal S}(D^2).
\end{equation}
Replacing the supremum norms by the John norms changes the degree by
$O_{p,\mathcal S}(D\log(D+2))$.  Changing the fixed adelic metric or the fixed
saturated integral model on $\calF_{\mathcal S}$ changes it by
$O_{p,\mathcal S}(D)$.
\end{proposition}

\begin{proof}
The curve $X_p$ is geometrically integral and normal, the line
$L=\mathcal O(P)$ is ample, and every local metric of
$\overline L_\alpha$ is a dominated semipositive metric: at finite places
it is a model metric, while at infinity it is the continuous Green metric of
\cref{prop:exact-mixed-divisor}.  The vector-bundle arithmetic
Hilbert--Samuel theorem
\cite[Corollary~4.5.2]{ChenMoriwaki2024} therefore applies directly to the
adelic vector bundle $\overline{\calF}_{\mathcal S}$ and gives
\[
 \lim_{D\to\infty}
 \frac{\widehatdeg H^0(X_p,L^D\otimes\calF_{\mathcal S})}{D^2/2}
 =\rank(\calF_{\mathcal S})
   (\overline L_\alpha^{\,2})
 =m\Energy_\alpha.
\]
This is \eqref{eq:appendix-exact-source} for the adelic supremum norms.

If $r_D=\dim H^0(X_p,L^D\otimes\calF_{\mathcal S})$, John's theorem for a
symmetric convex body gives a determinant distortion at most
$\tfrac12r_D\log r_D$ at each Archimedean place.  Since $r_D=mD+O(1)$, this
is $O(D\log D)$.  Two fixed dominated metrics on
$\calF_{\mathcal S}$ are uniformly equivalent place by place and agree
integrally outside a fixed finite set; their determinant distortion is
$O(r_D)=O(D)$.  The same applies to two fixed saturated models of the
coefficient bundle.  This proves all assertions.
\end{proof}

\begin{remark}[Determinant versus theta asymptotics]
The input used here is the determinant-degree formula of
\cite[Corollary~4.5.2]{ChenMoriwaki2024}.  The arithmetic Hilbert--Samuel
formula in \cite[Theorem~8.2.5]{BostCharles2022} is instead a formula for a
John-theta invariant.  No unproved identification of that theta invariant
with the determinant degree of the present source lattice is used.
\end{remark}

\subsection{Uniform finite-order jets and the \texorpdfstring{$\Theta$}{Theta}-cancellation}

Let $\overline C_{\mathcal S}$ denote the fixed metrized line obtained from
$\calF_{\mathcal S}|_P$ and from the fixed contraction pairing.  For a pivot
$n$, the natural target is
\begin{equation}\label{eq:natural-jet-line-app}
 \overline J_{D,n}^\alpha
 =P^*\overline L_\alpha^{\,D}
  \otimes\overline T_\alpha^{\vee\otimes n}
  \otimes\overline C_{\mathcal S}.
\end{equation}
The source quotient is endowed with the quotient norm from
$\overline M_D^\alpha$.  Choose any nonzero rational vector in
$J_{D,n}^\alpha$ and use it to identify the target with the standard scalar
line.  Denote the resulting local and global heights by
$h_{v,\alpha}^{\mathrm{sc}}$ and $h_\alpha^{\mathrm{sc}}$.  The product
formula gives the exact, trivialization-independent identity
\begin{equation}\label{eq:scalarization-identity}
 h_\alpha^{\mathrm{sc}}(\psi_D^{(n)})
 =h_\alpha^{\mathrm{nat}}(\psi_D^{(n)})
  +\widehatdeg\overline J_{D,n}^\alpha.
\end{equation}

At the distinguished place $p$ set $\delta_{p,D,n}=0$; its analytic
operator norm is controlled directly by the capacitary metric.  At every
finite place $v\ne p$, retain the coefficient excess
$\delta_{v,D,n}$ defined in \cref{sec:curve-bost}.

\begin{lemma}[Exact finite-place capacitary Schwarz estimate]
\label{lem:finite-order-cap-exact}
At the distinguished continuation place $p$, uniformly for all $D\ge0$ and
all $n\ge0$,
\begin{equation}\label{eq:finite-order-cap-local}
 \log\|\psi_D^{(n)}\|_{p,\mathrm{nat}}
 \le O_{p,\mathcal S}(1).
\end{equation}
At every finite place $v\ne p$,
\begin{equation}\label{eq:finite-order-coeff-local}
 \log\|\psi_D^{(n)}\|_{v,\mathrm{nat}}
 \le\delta_{v,D,n}+O_{v,p,\mathcal S}(1),
\end{equation}
and outside a fixed finite set the error term is zero.
\end{lemma}

\begin{proof}
At $p$ use the nested rational lemniscate
\eqref{eq:nested-lemniscate}.  If $1_P$ is the canonical rational section of
$\mathcal O(P)$ and the polar divisor of $u_p$ is $a_pP$, then
\begin{equation}\label{eq:unit-lemniscate-frame}
 \epsilon_p=u_p1_P^{a_p}
\end{equation}
is nonvanishing on $\Omega_p^\circ$ and has capacitary norm identically one
there.  Indeed, \cite[Corollary~5.8]{BostChambertLoir2009} gives
$\|1_P\|_p=|u_p|_p^{-1/a_p}$ on $\{|u_p|_p>1\}$.

Let a source vector have norm at most one and contract it with the horizontal
leaf.  Write $D=a_pk+b$, $0\le b<a_p$.  Dividing by $\epsilon_p^k$
trivializes the growing power of $\mathcal O(P)$; the residual power
$\mathcal O(bP)$, the coefficient bundle, and the contraction line are fixed.
For $r>1$ in the divisible value group, the maximum-principle computation in
the proof of \cite[Proposition~5.15]{BostChambertLoir2009}, applied on the
affinoid $\{|u_p|_p\ge r\}$, gives
\[
 \log\|\psi_D^{(n)}\|_{p,\mathrm{nat}}
 \le C_{p,\mathcal S}(r)+\frac{n}{a_p}\log r.
\]
By \eqref{eq:nested-lemniscate}, all these affinoids for $r>1$ sufficiently
close to one lie in the fixed affinoid
$\{|u_p|_p\ge1\}\Subset\Omega_p^+$.  The horizontal leaf and the fixed
contraction pairing are bounded on a neighborhood of that affinoid, while the
global source supremum norm bounds the coefficient section there.  Hence
$C_{p,\mathcal S}(r)$ is bounded above by one constant independent of
$r,D,n$.  Taking the infimum as $r\downarrow1$ removes the last term and
proves \eqref{eq:finite-order-cap-local}.  This is the finite-order statement
needed here; it is stronger than merely identifying the limiting canonical
tangent norm.

For $v\ne p$, the continuation domain is the integral Tate residue disc.  By
the local model prescription in the construction of $\mathscr X_\alpha$, the
finite metrics on $P^*\overline L_\alpha$ and $\overline T_\alpha$ are the
same integral target metrics used to define $\delta_{v,D,n}$ in
\cref{sec:curve-bost}.  Thus
\eqref{eq:finite-order-coeff-local} is the definition of the positive excess,
up to the bounded comparison of the fixed coefficient bundle and contraction
line.  These comparisons are integral outside a fixed finite set.
\end{proof}

\begin{lemma}[Archimedean Green--Jensen estimate]
\label{lem:arch-exact-jet}
Uniformly for every divided jet occurring in the vanishing filtration,
\begin{equation}\label{eq:arch-exact-jet}
 \sum_{\sigma}
 \log\|\psi_D^{(n)}\|_{\sigma,\mathrm{nat}}
 \le D\mathscr W_\alpha+O_{p,\mathcal S}(1).
\end{equation}
\end{lemma}

\begin{proof}
Contract first with the fixed horizontal leaf.  On $\Omega_\sigma$ this gives
a scalar analytic section $s$ of $\iota_\sigma^*L_\alpha^D$ tensored with a
fixed line.  Suppose that its first nonzero divided jet at $o_\sigma$ has
order $n$.  Poincar\'e--Lelong and Green's identity for $g_\sigma$ give
\begin{align*}
 \log\|j_{o_\sigma}^ns\|_{\mathrm{nat}}
 &\le
 \int_{\partial\Omega_\sigma}\log\|s\|\,\dd\mu_\sigma
 +D\int_{\Omega_\sigma}
   g_\sigma\,\iota_\sigma^*c_1(\overline L_{\alpha,\sigma})
 +O_{p,\mathcal S}(1).
\end{align*}
Every zero away from $o_\sigma$ occurs with the favorable sign and has been
discarded.  The source capacitary tangent metric in
\eqref{eq:natural-jet-line-app} is exactly the normalization that removes the
local derivative term.  This is the intrinsic Green--Jensen calculation
underlying \cite[Lemma~7.4.1]{CDT2024}.

By \cref{prop:positive-width}, the image of the closed pointed domain lies in
a slightly larger continuation domain.  The horizontal leaf and the fixed
coefficient-bundle pairing are therefore bounded on a neighborhood of the
boundary.  For a source vector of John norm at most one, its supremum norm is
at most one because the John ellipsoid is contained in the supremum unit
ball.  The boundary integral is consequently $O_{p,\mathcal S}(1)$,
uniformly in $D$ and $n$.  Summing over $\sigma$ and using the projection
formula in \eqref{eq:energy-theta} proves \eqref{eq:arch-exact-jet}.
\end{proof}

\begin{proposition}[Uniform single-pivot cancellation]
\label{prop:exact-single-jet}
Fix $C>0$.  Uniformly for pivots $0\le n\le CD$,
\begin{align}
 \widehatdeg\overline J_{D,n}^\alpha
 &=D\Theta_\alpha-n\LambdaAH_\alpha+O_{p,\mathcal S}(1),
 \label{eq:app-target-degree}\\
 h_\alpha^{\mathrm{nat}}(\psi_D^{(n)})
 &\le D(\Energy_\alpha-\Theta_\alpha)+\delta_{D,n}
      +O_{C,p,\mathcal S}(1),
 \label{eq:app-natural-height}\\
 h_\alpha^{\mathrm{sc}}(\psi_D^{(n)})
 &\le D\Energy_\alpha-n\LambdaAH_\alpha+\delta_{D,n}
      +O_{C,p,\mathcal S}(1),
 \label{eq:app-scalar-height}\\
 \widehatdeg\overline G_{D,n}^\alpha
 &\le D\Energy_\alpha-n\LambdaAH_\alpha+\delta_{D,n}
      +O_{C,p,\mathcal S}(1).
 \label{eq:app-graded-degree}
\end{align}
Thus the $D\Theta_\alpha$ target contribution cancels exactly with the
$-D\Theta_\alpha$ part of the natural evaluation estimate.
\end{proposition}

\begin{proof}
The target-degree identity follows immediately from
\eqref{eq:natural-jet-line-app}.  Summing
\cref{lem:finite-order-cap-exact} over the finite places contributes
$\delta_{D,n}+O_{p,\mathcal S}(1)$.  Adding
\cref{lem:arch-exact-jet} and using
$\mathscr W_\alpha=\Energy_\alpha-\Theta_\alpha$ proves
\eqref{eq:app-natural-height}.

Combining \eqref{eq:scalarization-identity},
\eqref{eq:app-target-degree}, and \eqref{eq:app-natural-height} cancels
$D\Theta_\alpha$ and proves \eqref{eq:app-scalar-height}.  For the scalarized
nonzero map from the rank-one source quotient to the standard degree-zero
line, the rank-one degree identity is
\[
 \widehatdeg\overline G_{D,n}^\alpha
 =h_\alpha^{\mathrm{sc}}(\psi_D^{(n)}).
\]
This proves \eqref{eq:app-graded-degree}.
\end{proof}

\begin{remark}[Where overflow is counted]
The finite Bost--Chambert-Loir vertical corrections are contained in the target
degree $D\Theta_\alpha$.  The Archimedean direct-image contribution enters
$D\mathscr W_\alpha=D(\Energy_\alpha-\Theta_\alpha)$.  Their sum is
$D\Energy_\alpha$.  No formal or Archimedean term is inserted a second time.
This is the precise target-normal cancellation asserted in
\cref{thm:curve-Bost-exact}.
\end{remark}

\subsection{Hybrid lattices and exact slopes}

\begin{proposition}[Exact hybrid assembly]
\label{prop:exact-hybrid-assembly}
For the metrics above and the hybrid lattice
\eqref{eq:curve-hybrid-lattice},
\begin{equation}\label{eq:app-exact-hybrid}
 \widehatdeg\overline M_D^{\alpha,\mathrm{hyb}}
 =\left(\frac m2\Energy_\alpha-\xi Q_{\mathcal S}\right)D^2
  +o_{p,\mathcal S}(D^2).
\end{equation}
Moreover
\begin{equation}\label{eq:app-exact-degree-sum}
 \sum_{n\in V_D}\widehatdeg\overline G_{D,n}^\alpha
 \le\left(-\frac{m^2}{2}\LambdaAH_\alpha
          +m\Energy_\alpha\right)D^2
    +\mathscr A_D+o_{p,\mathcal S}(D^2),
\end{equation}
and joint rationality implies
\begin{equation}\label{eq:app-exact-slopes}
 \tau^{\mathrm{curve}}_{p,\mathcal S}(\xi)
 \ge\LambdaAH_\alpha-\frac{\Energy_\alpha}{m}.
\end{equation}
\end{proposition}

\begin{proof}
Use the finite $\Q$-model of $\overline L_\alpha$ as the reference source
model.  At every place $\ell\ne p$, choose the coefficient-bundle models in
\cref{sec:qjet} from this same model, after enlarging the fixed exceptional
set.  Only the fixed metric and integral model of
$\calF_{\mathcal S}$ then require comparison, which changes determinant
degrees and summed fixed-place heights by
$O_{p,\mathcal S}(D\log(D+2))$.  At the distinguished place $p$ no comparison
with a coefficient lattice is made: it is controlled by the capacitary
estimate and is excluded from $\delta_{D,n}$.  Passing to a sublattice can
only improve its analytic operator norm.

The positive block of weight $s$ has rank $sD+O_{p,s}(1)$.  Replacing it by
$L_N^sM_{D,s}^+$ therefore changes the degree by
\[
 -s^2D\log L_N+O_{p,s}(\log L_N).
\]
Since $\log L_N=\xi D+o(D)$, summing over $s$ and using
\cref{prop:exact-source-formula} proves \eqref{eq:app-exact-hybrid}.
Shrinking the finite source lattice increases its gauge norm, so the natural
jet estimate of \cref{prop:exact-single-jet} remains valid.

The residual small-prime heights are $o(D^2)$ by the proper-prime-power
argument in \cref{prop:curve-hybrid-compatibility}.  At primes
$\ell>\xi D$, one has $\ell\ne p$ for $D\gg1$, and the compatible coefficient
models permit the literal application of
\cref{prop:q-jet-Cartier,prop:curve-large}; hence
\[
 \sum_{n\in V_D}\delta_{D,n}
 \le s_*mG_{d_p,m}(\xi)D^2+o_{p,\mathcal S}(D^2).
\]
By \cref{thm:finite-map-zero}, all pivots satisfy $n\le CD$ for one fixed
$C=C_{p,\mathcal S}$.  Sum \eqref{eq:app-graded-degree} over the
$mD+O(1)$ pivots.  Its uniform $O_{C,p,\mathcal S}(1)$ remainder contributes
only $O_{p,\mathcal S}(D)$.  Together with
\eqref{eq:robust-pivot-sum}, this proves
\eqref{eq:app-exact-degree-sum}.

The vanishing filtration has saturated rank-one quotients whose degrees add:
\[
 \widehatdeg\overline M_D^{\alpha,\mathrm{hyb}}
 =\sum_{n\in V_D}\widehatdeg\overline G_{D,n}^\alpha.
\]
Comparing \eqref{eq:app-exact-hybrid} with
\eqref{eq:app-exact-degree-sum} gives
\[
 \frac m2\Energy_\alpha-\xi Q_{\mathcal S}
 \le m\Energy_\alpha-\frac{m^2}{2}\LambdaAH_\alpha
    +s_*mG_{d_p,m}(\xi).
\]
Rearranging and using \eqref{eq:curve-tau} proves
\eqref{eq:app-exact-slopes}.
\end{proof}

\begin{remark}[Logical role]
The exact theorem is a canonical refinement of the general-curve slopes
argument.  The proof of \cref{thm:density-all-primes} remains based on the
robust inequality \eqref{eq:robust-curve-slopes}; hence no density conclusion
depends on the more delicate exact-normalization appendix.
\end{remark}

\section{Derivation of the capped window}\label{app:capped}

For completeness, let $m\ge2$ and $1<\xi<m$.  A prime
$\ell=xD$ contributes when a pivot lies in a degree-$D$ interval ending at a
positive multiple of $\ell$.  Maximal packing of distinct pivots gives
\[
 w_m(x)=K_x+\positivepart{m-(K_x+1)x},
 \qquad K_x=\floor{\frac{m-1}{x}}.
\]
Integrating $w_m$ gives \eqref{eq:I}.  The function is decreasing and crosses
the integer level $\rho$ at $x=m/(\rho+1)$.  Therefore
\[
 \int_\xi^m\min\{w_m(x),\rho\}\,\dd x
 =\rho\left(\frac{m}{\rho+1}-\xi\right)
 +I^m_{m/(\rho+1)}\left(\frac{m}{\rho+1}\right)
\]
when $\xi<m/(\rho+1)$.  At the crossing,
$K=\rho$ and the final quadratic term vanishes, giving
\[
 I^m_{m/(\rho+1)}
 =\left(m-\frac12\right)H_\rho-\frac{m\rho}{\rho+1}.
\]
This proves \eqref{eq:I-cap-explicit}.

\section{Certificate methodology}\label{app:certificate}

Real intervals are represented by integer endpoints at scale $10^{42}$ and
every operation is rounded outward.  The constants are enclosed from
\[
 \pi=16\arctan\frac15-4\arctan\frac1{239},
 \qquad
 \log p=2\operatorname{arctanh}\frac{p-1}{p+1}.
\]
The alternating arctangent remainder is bounded by the first omitted term; the
positive arctanh tail is bounded geometrically.  For rational $x\in(0,1)$,
\[
 \arccos x=2\arctan\sqrt{\frac{1-x}{1+x}},
\]
and square roots are enclosed by exact integer square root.  Harmonic sums and
all algebraic terms are rational.  The active collision layers are selected
by the exact test $cY<1$.

For fixed $p$, stage $k$, largest candidate $x$, and previous-entry sum $T$,
\cref{lem:convex-search} gives the exact maximal $Q$.  The program evaluates
the upper endpoint of the arithmetic cost and the lower endpoint of the
analytic margin.  A strictly positive lower endpoint excludes the entire
aggregate state.  The stage witnesses and every worst state are printed in
the stored output.

\begin{thebibliography}{99}

\bibitem{Bost1999}
J.-B.~Bost,
\newblock Potential theory and Lefschetz theorems for arithmetic surfaces,
\newblock \emph{Ann. Sci. \`Ecole Norm. Sup. (4)} \textbf{32} (1999), no.~2,
241--312.

\bibitem{Bost2001}
J.-B.~Bost,
\newblock Algebraic leaves of algebraic foliations over number fields,
\newblock \emph{Publ. Math. Inst. Hautes \`Etudes Sci.} \textbf{93} (2001),
161--221.


\bibitem{BostChambertLoir2009}
J.-B.~Bost and A.~Chambert-Loir,
\newblock Analytic curves in algebraic varieties over number fields,
\newblock in \emph{Algebra, Arithmetic, and Geometry: In Honor of
Yu.~I.~Manin}, vol.~I, Progress in Mathematics \textbf{269}, Birkh\"auser,
2009; arXiv:math/0702593.

\bibitem{BostCharles2022}
J.-B.~Bost and F.~Charles,
\newblock \emph{Quasi-Projective and Formal-Analytic Arithmetic Surfaces},
\newblock Annals of Mathematics Studies \textbf{224}, Princeton University
Press, 2026; arXiv:2206.14242v2.

\bibitem{Buzzard2003}
K.~Buzzard,
\newblock Analytic continuation of overconvergent eigenforms,
\newblock \emph{J. Amer. Math. Soc.} \textbf{16} (2003), 29--55.

\bibitem{Calegari2005}
F.~Calegari,
\newblock Irrationality of certain $p$-adic periods for small $p$,
\newblock \emph{Int. Math. Res. Not.} 2005, no.~20, 1235--1249.

\bibitem{CDT2024}
F.~Calegari, V.~Dimitrov and Y.~Tang,
\newblock The linear independence of $1$, $\zeta(2)$, and $L(2,\chi_{-3})$,
\newblock arXiv:2408.15403v2 (2024).

\bibitem{CDT2025}
F.~Calegari, V.~Dimitrov and Y.~Tang,
\newblock Arithmetic holonomy bounds and effective Diophantine approximation,
\newblock arXiv:2510.04156 (2025).

\bibitem{ChenMoriwaki2024}
H.~Chen and A.~Moriwaki,
\newblock \emph{Positivity in Arakelov Geometry over Adelic Curves:
Hilbert--Samuel Formula and Equidistribution Theorem},
\newblock Progress in Mathematics \textbf{355}, Birkh\"auser, 2024;
see also arXiv:2207.02033, Corollary~4.5.2.

\bibitem{GPRJ2025}
A.~Graham, V.~Pilloni and J.~Rodrigues Jacinto,
\newblock $p$-adic interpolation of Gauss--Manin connections on nearly
overconvergent modular forms and $p$-adic $L$-functions,
\newblock \emph{Compos. Math.} \textbf{161} (2025), 2380--2441.

\bibitem{Katz1973}
N.~M.~Katz,
\newblock $p$-adic properties of modular schemes and modular forms,
\newblock in \emph{Modular Functions of One Variable III}, Lecture Notes in
Math. \textbf{350}, Springer, 1973, pp.~69--190.


\bibitem{Cesnavicius2017}
K.~\v{C}esnavi\v{c}ius,
\newblock A modular description of $\mathscr X_0(n)$,
\newblock \emph{Algebra Number Theory} \textbf{11} (2017), no.~9, 2001--2089;
\newblock arXiv:1511.07475.

\bibitem{ColemanMcMurdy2006}
R.~Coleman and K.~McMurdy,
\newblock Fake CM and the stable model of $X_0(Np^3)$,
\newblock \emph{Doc. Math.} 2006, Extra Volume, 261--300; in particular the
review of the supersingular annuli of $X_0(p)$ in \S2.

\bibitem{Deuring1941}
M.~Deuring,
\newblock Die Typen der Multiplikatorenringe elliptischer Funktionenk\"orper,
\newblock \emph{Abh. Math. Sem. Univ. Hamburg} \textbf{14} (1941), 197--272.

\bibitem{Maier2009}
R.~S.~Maier,
\newblock On rationally parametrized modular equations,
\newblock \emph{J. Ramanujan Math. Soc.} \textbf{24} (2009), no.~1, 1--73;
\newblock arXiv:math/0611041.

\bibitem{RodriguezCamargo2023}
J.~E.~Rodr\'iguez Camargo,
\newblock $p$-adic Eichler--Shimura maps for the modular curve,
\newblock \emph{Compos. Math.} \textbf{159} (2023), no.~6, 1214--1249.

\bibitem{Urban2013}
E.~Urban,
\newblock Nearly overconvergent modular forms,
\newblock in \emph{Iwasawa Theory 2012}, Contributions in Mathematical and
Computational Sciences \textbf{7}, Springer, 2014, pp.~401--441.

\end{thebibliography}

\end{document}
